\documentclass[a4paper,10pt]{article}
\usepackage{amsmath,amssymb}
\usepackage{fontspec}
\setmainfont{Liberation Serif}
\usepackage{graphicx}
\usepackage{eso-pic}
\usepackage{tikz}
\usepackage[most]{tcolorbox}
\usepackage{geometry}
\usepackage{fancyhdr}
\usepackage{lastpage}
\usepackage{booktabs}
\usepackage{listings}
\usepackage{enumitem}
\usepackage{microtype}

\geometry{
    a4paper,
    top=1.5cm,
    bottom=1.8cm,
    left=1.2cm,
    right=1.2cm
}

% Persistent Unremovable Canvas Watermark on Every Page (Shared XObject)
\newsavebox{\watermarkbox}
\sbox{\watermarkbox}{%
    \begin{tabular}{c}
        \includegraphics[width=8cm]{/home/johan/GATE/GATE_CS_PYQ_DB/assets/studiva_butterfly_icon_transparent.png} \\[-0.5em]
        {\fontsize{18}{22}\selectfont\bfseries\scshape www.studiva.co.in}
    \end{tabular}%
}
\AddToShipoutPictureBG{%
    \begin{tikzpicture}[remember picture, overlay]
        \node[opacity=0.12] at (current page.center) {\usebox{\watermarkbox}};
    \end{tikzpicture}%
}

% Header and Footer
\pagestyle{fancy}
\fancyhf{}
\fancyhead[L]{\small\scshape GATE 2025 Computer Science and Information Technology (CS2)}
\fancyhead[R]{\small\scshape Practice Question Bank with Solutions}
\fancyfoot[L]{\scriptsize\textsc{Copyright \copyright\ 2026} Studiva\textsuperscript{TM}}
\fancyfoot[C]{\small Page \thepage\ of \pageref{LastPage}}
\fancyfoot[R]{\raisebox{-0.2em}{\includegraphics[height=1.0em]{/home/johan/GATE/GATE_CS_PYQ_DB/assets/studiva_butterfly_icon_transparent.png}}\hspace{0.35em}\scriptsize\texttt{www.studiva.co.in}}
\renewcommand{\headrulewidth}{0.4pt}
\renewcommand{\footrulewidth}{0.4pt}

% Code Listing Style (Monochrome, clean)
\lstset{
    basicstyle=\footnotesize\ttfamily,
    numbers=none,
    frame=none,
    tabsize=2,
    breaklines=true,
    showstringspaces=false
}

\begin{document}

% ==================== PAGE 1: DEDICATED IMPRINT & COPYRIGHT PAGE ====================
\thispagestyle{empty}
\vspace*{0.5cm}
\begin{center}
    \includegraphics[width=2.5cm]{/home/johan/GATE/GATE_CS_PYQ_DB/assets/studiva_butterfly_icon_transparent.png}\\[0.6em]
    {\LARGE \textbf{GATE 2025 COMPUTER SCIENCE \& IT}}\\[0.3em]
    {\Large \textbf{PREVIOUS YEAR QUESTION BANK WITH SOLUTIONS}}\\[0.2em]
    {\normalsize \textsc{Session CS2 $\cdot$ Examination Year 2025}}\\[1.2em]
    \rule{0.6\textwidth}{0.6pt}
\end{center}

\vspace{1em}
\noindent\begin{tabular}{ll}
\textbf{Author / Compiler:} & Rupam Ghosh \\
\textbf{Publisher \& Platform:} & Studiva\textsuperscript{TM} (\texttt{www.studiva.co.in}) \\
\textbf{Official Inquiries:} & \texttt{contact@studiva.co.in} \\
\textbf{Edition:} & First Standard Edition (2026) \\
\textbf{Syllabus Reference:} & Official GATE 2026 Syllabus (Organizing Institute: IIT Guwahati) \\
\end{tabular}

\vspace{1.5em}
\begin{tcolorbox}[enhanced, standard jigsaw, opacityback=0, colframe=black!40, arc=1mm, boxrule=0.5pt, left=10pt, right=10pt, top=8pt, bottom=8pt]
\textbf{\large COPYRIGHT \& INTELLECTUAL PROPERTY NOTICE}\\[0.6em]
Copyright \copyright\ 2026 Rupam Ghosh \& Studiva\textsuperscript{TM}. All Rights Reserved.\\[0.5em]
This educational compilation, problem organization, taxonomy classification, LaTeX typesetting, hints, and explanatory solutions are the proprietary intellectual property of \textbf{Rupam Ghosh} and \textbf{Studiva\textsuperscript{TM}}.\\[0.6em]
\textbf{Terms of Use \& Attribution Rules:}
\begin{itemize}[leftmargin=*, nosep]
    \item \textbf{Personal Educational Use:} Permitted freely for individual students and educators preparing for competitive examinations.
    \item \textbf{Mandatory Attribution:} Any reproduction, republication, hosting, or sharing (digital or physical) must prominently credit \textbf{Rupam Ghosh \& Studiva\textsuperscript{TM}} and include an active, direct link to \texttt{https://www.studiva.co.in}.
    \item \textbf{Commercial Restrictions:} Commercial resale, sublicensing, course bundling, or inclusion in paid test series without prior written permission is strictly prohibited.
    \item \textbf{Anti-Tampering Notice:} Removing, altering, or obscuring copyright notices, page watermarks, or website links from this document is strictly prohibited.
\end{itemize}
\vspace{0.6em}
\textbf{Fair Use Disclaimer:}\\
The underlying examination questions originate from the Graduate Aptitude Test in Engineering (GATE), conducted by the Indian Institutes of Technology (IITs) and Indian Institute of Science (IISc). Question texts are reproduced for non-commercial educational study under Fair Use doctrines. No copyright is claimed over the original GATE examination question texts.
\end{tcolorbox}

\vspace{1.2em}
\begin{center}
\small\itshape
Practice Recommendation: Use an opaque sheet of paper or ruler to cover the right-hand column during mock attempts. Tick your chosen answers directly in the check boxes on the left, then slide the cover to check the solution and hint.
\end{center}

\newpage
% ==================== PAGE 2 ONWARDS: QUESTION BANK WITH SOLUTIONS ====================

\vspace{6pt}
\subsection*{General Aptitude}
\vspace{-4pt}
\noindent
\begin{tcolorbox}[
    blanker,
    sidebyside,
    sidebyside align=top,
    lefthand width=0.55\textwidth,
    top=4pt, bottom=4pt, left=0pt, right=0pt,
    middle=16pt,
    underlay={
        \draw[dashed, black!80, line width=0.8pt, dash pattern=on 4pt off 3pt]
            ([xshift=0.55\textwidth+8pt]frame.north west) --
            ([xshift=0.55\textwidth+8pt]frame.south west);
    }
]
\textbf{Q.1} \quad {\footnotesize\textsc{[MCQ $\cdot$ 1 Mark(s) $\cdot$ Neg: 0.33]}} \hfill {\tiny\texttt{2025\_CS2\_GA\_Q01}}
\vspace{2pt}\hrule height 0.3pt\vspace{4pt}

Despite his initial hesitation, Rehman’s \_\_\_\_\_\_\_\_\_ to contribute to the success of the project never wavered.
Select the most appropriate option to complete the above sentence.

\par\vspace{6pt}\noindent
\begin{tabular}{@{}p{0.48\linewidth}@{\hspace{0.04\linewidth}}p{0.48\linewidth}@{}}
$\square$\ \textbf{(A)} ambivalence & $\square$\ \textbf{(B)} satisfaction \\
\addlinespace[3pt]
$\square$\ \textbf{(C)} resolve & $\square$\ \textbf{(D)} revolve
\end{tabular}

\tcblower

\textbf{Official Key:} \textbf{\large C} \hfill \textit{\small [Easy]}
\vspace{2pt}\hrule height 0.3pt\vspace{3pt}
\textbf{Topic:} Verbal Aptitude \\
\textbf{Subtopic:} \textit{Vocabulary: words, idioms, and phrases in context}
\vspace{2pt}\hrule height 0.3pt\vspace{3pt}
\textbf{Guide / Hint:} \\
{\small The sentence sets up a contrast between Rehman's initial hesitation and his continuing commitment. 'Resolve' denotes firm determination and perseverance, which fits the context accurately.}

\end{tcolorbox}
\vspace{6pt}\hrule height 0.35pt\vspace{6pt}
\noindent
\begin{tcolorbox}[
    blanker,
    sidebyside,
    sidebyside align=top,
    lefthand width=0.55\textwidth,
    top=4pt, bottom=4pt, left=0pt, right=0pt,
    middle=16pt,
    underlay={
        \draw[dashed, black!80, line width=0.8pt, dash pattern=on 4pt off 3pt]
            ([xshift=0.55\textwidth+8pt]frame.north west) --
            ([xshift=0.55\textwidth+8pt]frame.south west);
    }
]
\textbf{Q.2} \quad {\footnotesize\textsc{[MCQ $\cdot$ 1 Mark(s) $\cdot$ Neg: 0.33]}} \hfill {\tiny\texttt{2025\_CS2\_GA\_Q02}}
\vspace{2pt}\hrule height 0.3pt\vspace{4pt}

$\text{Bird} : \text{Nest} :: \text{Bee} : \text{\_\_\_\_\_\_\_}$
Select the correct option to complete the analogy.

\par\vspace{6pt}\noindent
\begin{tabular}{@{}p{0.48\linewidth}@{\hspace{0.04\linewidth}}p{0.48\linewidth}@{}}
$\square$\ \textbf{(A)} Kennel & $\square$\ \textbf{(B)} Hammock \\
\addlinespace[3pt]
$\square$\ \textbf{(C)} Hive & $\square$\ \textbf{(D)} Lair
\end{tabular}

\tcblower

\textbf{Official Key:} \textbf{\large C} \hfill \textit{\small [Easy]}
\vspace{2pt}\hrule height 0.3pt\vspace{3pt}
\textbf{Topic:} Verbal Aptitude \\
\textbf{Subtopic:} \textit{Vocabulary: word analogy and contextual relations}
\vspace{2pt}\hrule height 0.3pt\vspace{3pt}
\textbf{Guide / Hint:} \\
{\small This is a habitat/dwelling analogy: a bird naturally lives in a nest, and similarly, a bee lives in a hive.}

\end{tcolorbox}
\vspace{6pt}\hrule height 0.35pt\vspace{6pt}
\noindent
\begin{tcolorbox}[
    blanker,
    sidebyside,
    sidebyside align=top,
    lefthand width=0.55\textwidth,
    top=4pt, bottom=4pt, left=0pt, right=0pt,
    middle=16pt,
    underlay={
        \draw[dashed, black!80, line width=0.8pt, dash pattern=on 4pt off 3pt]
            ([xshift=0.55\textwidth+8pt]frame.north west) --
            ([xshift=0.55\textwidth+8pt]frame.south west);
    }
]
\textbf{Q.3} \quad {\footnotesize\textsc{[MCQ $\cdot$ 1 Mark(s) $\cdot$ Neg: 0.33]}} \hfill {\tiny\texttt{2025\_CS2\_GA\_Q03}}
\vspace{2pt}\hrule height 0.3pt\vspace{4pt}

If $P e^x = Q e^{-x}$ for all real values of $x$, which one of the following statements is true?

\par\vspace{6pt}\noindent
\begin{tabular}{@{}p{0.48\linewidth}@{\hspace{0.04\linewidth}}p{0.48\linewidth}@{}}
$\square$\ \textbf{(A)} $P = Q = 0$ & $\square$\ \textbf{(B)} $P = Q = 1$ \\
\addlinespace[3pt]
$\square$\ \textbf{(C)} $P = 1; Q = -1$ & $\square$\ \textbf{(D)} $\frac{P}{Q} = 0$
\end{tabular}

\tcblower

\textbf{Official Key:} \textbf{\large A} \hfill \textit{\small [Medium]}
\vspace{2pt}\hrule height 0.3pt\vspace{3pt}
\textbf{Topic:} Quantitative Aptitude \\
\textbf{Subtopic:} \textit{Numerical computation and functions: exponents and equations}
\vspace{2pt}\hrule height 0.3pt\vspace{3pt}
\textbf{Guide / Hint:} \\
{\small For $P e^x = Q e^{-x}$ to hold for all real $x$, multiply both sides by $e^x$ to get $P e^{2x} = Q$. Differentiating with respect to $x$ gives $2P e^{2x} = 0$, implying $P = 0$. Substituting $P = 0$ yields $Q = 0$, so $P = Q = 0$.}

\end{tcolorbox}
\vspace{6pt}\hrule height 0.35pt\vspace{6pt}
\noindent
\begin{tcolorbox}[
    blanker,
    sidebyside,
    sidebyside align=top,
    lefthand width=0.55\textwidth,
    top=4pt, bottom=4pt, left=0pt, right=0pt,
    middle=16pt,
    underlay={
        \draw[dashed, black!80, line width=0.8pt, dash pattern=on 4pt off 3pt]
            ([xshift=0.55\textwidth+8pt]frame.north west) --
            ([xshift=0.55\textwidth+8pt]frame.south west);
    }
]
\textbf{Q.5} \quad {\footnotesize\textsc{[MCQ $\cdot$ 1 Mark(s) $\cdot$ Neg: 0.33]}} \hfill {\tiny\texttt{2025\_CS2\_GA\_Q05}}
\vspace{2pt}\hrule height 0.3pt\vspace{4pt}

Let $p_1$ and $p_2$ denote two arbitrary prime numbers. Which one of the following statements is correct for all values of $p_1$ and $p_2$?

\par\vspace{6pt}\noindent
\begin{tabular}{@{}p{0.98\linewidth}@{}}
$\square$\ \textbf{(A)} $p_1 + p_2$ is not a prime number. \\
\addlinespace[3pt]
$\square$\ \textbf{(B)} $p_1 p_2$ is not a prime number. \\
\addlinespace[3pt]
$\square$\ \textbf{(C)} $p_1 + p_2 + 1$ is a prime number. \\
\addlinespace[3pt]
$\square$\ \textbf{(D)} $p_1 p_2 + 1$ is a prime number.
\end{tabular}

\tcblower

\textbf{Official Key:} \textbf{\large B} \hfill \textit{\small [Easy]}
\vspace{2pt}\hrule height 0.3pt\vspace{3pt}
\textbf{Topic:} Quantitative Aptitude \\
\textbf{Subtopic:} \textit{Elementary number theory and prime numbers}
\vspace{2pt}\hrule height 0.3pt\vspace{3pt}
\textbf{Guide / Hint:} \\
{\small For any two prime numbers $p_1$ and $p_2$, the product $p_1 p_2$ has at least three distinct positive divisors ($1, p_1, p_1 p_2$ when $p_1=p_2$, or $1, p_1, p_2, p_1 p_2$ when $p_1 \neq p_2$). Hence, $p_1 p_2$ can never be prime.}

\end{tcolorbox}
\vspace{6pt}\hrule height 0.35pt\vspace{6pt}
\noindent
\begin{tcolorbox}[
    blanker,
    sidebyside,
    sidebyside align=top,
    lefthand width=0.55\textwidth,
    top=4pt, bottom=4pt, left=0pt, right=0pt,
    middle=16pt,
    underlay={
        \draw[dashed, black!80, line width=0.8pt, dash pattern=on 4pt off 3pt]
            ([xshift=0.55\textwidth+8pt]frame.north west) --
            ([xshift=0.55\textwidth+8pt]frame.south west);
    }
]
\textbf{Q.6} \quad {\footnotesize\textsc{[MCQ $\cdot$ 2 Mark(s) $\cdot$ Neg: 0.66]}} \hfill {\tiny\texttt{2025\_CS2\_GA\_Q06}}
\vspace{2pt}\hrule height 0.3pt\vspace{4pt}

Based only on the conversation below, identify the logically correct inference:
``Even if I had known that you were in the hospital, I would not have gone there to see you'', Ramya told Josephine.

\par\vspace{6pt}\noindent
\begin{tabular}{@{}p{0.98\linewidth}@{}}
$\square$\ \textbf{(A)} Ramya knew that Josephine was in the hospital. \\
\addlinespace[3pt]
$\square$\ \textbf{(B)} Ramya did not know that Josephine was in the hospital. \\
\addlinespace[3pt]
$\square$\ \textbf{(C)} Ramya and Josephine were once close friends; but now, they are not. \\
\addlinespace[3pt]
$\square$\ \textbf{(D)} Josephine was in the hospital due to an injury to her leg.
\end{tabular}

\tcblower

\textbf{Official Key:} \textbf{\large B} \hfill \textit{\small [Easy]}
\vspace{2pt}\hrule height 0.3pt\vspace{3pt}
\textbf{Topic:} Analytical Aptitude \\
\textbf{Subtopic:} \textit{Logic: deduction and conditional statements}
\vspace{2pt}\hrule height 0.3pt\vspace{3pt}
\textbf{Guide / Hint:} \\
{\small The subjunctive statement 'Even if I had known that you were in the hospital...' logically presupposes that the speaker did not know about Josephine's hospitalization at that time.}

\end{tcolorbox}
\vspace{6pt}\hrule height 0.35pt\vspace{6pt}
\noindent
\begin{tcolorbox}[
    blanker,
    sidebyside,
    sidebyside align=top,
    lefthand width=0.55\textwidth,
    top=4pt, bottom=4pt, left=0pt, right=0pt,
    middle=16pt,
    underlay={
        \draw[dashed, black!80, line width=0.8pt, dash pattern=on 4pt off 3pt]
            ([xshift=0.55\textwidth+8pt]frame.north west) --
            ([xshift=0.55\textwidth+8pt]frame.south west);
    }
]
\textbf{Q.7} \quad {\footnotesize\textsc{[MCQ $\cdot$ 2 Mark(s) $\cdot$ Neg: 0.66]}} \hfill {\tiny\texttt{2025\_CS2\_GA\_Q07}}
\vspace{2pt}\hrule height 0.3pt\vspace{4pt}

If $\text{IMAGE}$ and $\text{FIELD}$ are coded as $\text{FHBNJ}$ and $\text{EMFJG}$ respectively then, which one among the given options is the most appropriate code for $\text{BEACH}$?

\par\vspace{6pt}\noindent
\begin{tabular}{@{}p{0.48\linewidth}@{\hspace{0.04\linewidth}}p{0.48\linewidth}@{}}
$\square$\ \textbf{(A)} CEADP & $\square$\ \textbf{(B)} IDBFC \\
\addlinespace[3pt]
$\square$\ \textbf{(C)} JGIBC & $\square$\ \textbf{(D)} IBCEC
\end{tabular}

\tcblower

\textbf{Official Key:} \textbf{\large B} \hfill \textit{\small [Medium]}
\vspace{2pt}\hrule height 0.3pt\vspace{3pt}
\textbf{Topic:} Analytical Aptitude \\
\textbf{Subtopic:} \textit{Letter patterns, sequencing and coding-decoding}
\vspace{2pt}\hrule height 0.3pt\vspace{3pt}
\textbf{Guide / Hint:} \\
{\small The coding rule applies a $+1$ alphabetical shift to each character followed by string reversal. Shifting $\text{BEACH}$ by $+1$ gives $\text{CFBDI}$, and reversing this yields $\text{IDBFC}$.}

\end{tcolorbox}
\vspace{6pt}\hrule height 0.35pt\vspace{6pt}
\noindent
\begin{tcolorbox}[
    blanker,
    sidebyside,
    sidebyside align=top,
    lefthand width=0.55\textwidth,
    top=4pt, bottom=4pt, left=0pt, right=0pt,
    middle=16pt,
    underlay={
        \draw[dashed, black!80, line width=0.8pt, dash pattern=on 4pt off 3pt]
            ([xshift=0.55\textwidth+8pt]frame.north west) --
            ([xshift=0.55\textwidth+8pt]frame.south west);
    }
]
\textbf{Q.8} \quad {\footnotesize\textsc{[MCQ $\cdot$ 2 Mark(s) $\cdot$ Neg: 0.66]}} \hfill {\tiny\texttt{2025\_CS2\_GA\_Q08}}
\vspace{2pt}\hrule height 0.3pt\vspace{4pt}

Which one of the following options is correct for the given data in the table?
\par\medskip
\begin{center}\small
\begin{tabular}{cccc}
\toprule
Iteration ($i$) & Input ($I$) & Output ($X$) & Output ($Y$) \\
\midrule
0 & 20 & 20 & 20 \\
1 & $-4$ & 16 & $-80$ \\
2 & 10 & 26 & $-800$ \\
3 & 15 & 41 & $-12000$ \\
\bottomrule
\end{tabular}\end{center}

\par\vspace{6pt}\noindent
\begin{tabular}{@{}p{0.98\linewidth}@{}}
$\square$\ \textbf{(A)} $X(i) = X(i-1) + I(i); \; Y(i) = Y(i-1) I(i); \; i > 0$ \\
\addlinespace[3pt]
$\square$\ \textbf{(B)} $X(i) = X(i-1) I(i); \; Y(i) = Y(i-1) + I(i); \; i > 0$ \\
\addlinespace[3pt]
$\square$\ \textbf{(C)} $X(i) = X(i-1) I(i); \; Y(i) = Y(i-1) I(i); \; i > 0$ \\
\addlinespace[3pt]
$\square$\ \textbf{(D)} $X(i) = X(i-1) + I(i); \; Y(i) = Y(i-1) I(i-1); \; i > 0$
\end{tabular}

\tcblower

\textbf{Official Key:} \textbf{\large A} \hfill \textit{\small [Medium]}
\vspace{2pt}\hrule height 0.3pt\vspace{3pt}
\textbf{Topic:} Quantitative Aptitude \\
\textbf{Subtopic:} \textit{Numerical reasoning and recurrence relations}
\vspace{2pt}\hrule height 0.3pt\vspace{3pt}
\textbf{Guide / Hint:} \\
{\small Check the sequence updates: $X(i) = X(i-1) + I(i)$ yields $20 + (-4) = 16, 16 + 10 = 26, 26 + 15 = 41$, and $Y(i) = Y(i-1) \cdot I(i)$ yields $20 \times (-4) = -80, -80 \times 10 = -800, -800 \times 15 = -12000$.}

\end{tcolorbox}
\vspace{6pt}\hrule height 0.35pt\vspace{6pt}
\vspace{6pt}
\subsection*{Computer Science and Information Technology}
\vspace{-4pt}
\noindent
\begin{tcolorbox}[
    blanker,
    sidebyside,
    sidebyside align=top,
    lefthand width=0.55\textwidth,
    top=4pt, bottom=4pt, left=0pt, right=0pt,
    middle=16pt,
    underlay={
        \draw[dashed, black!80, line width=0.8pt, dash pattern=on 4pt off 3pt]
            ([xshift=0.55\textwidth+8pt]frame.north west) --
            ([xshift=0.55\textwidth+8pt]frame.south west);
    }
]
\textbf{Q.11} \quad {\footnotesize\textsc{[MCQ $\cdot$ 1 Mark(s) $\cdot$ Neg: 0.33]}} \hfill {\tiny\texttt{2025\_CS2\_CS\_Q11}}
\vspace{2pt}\hrule height 0.3pt\vspace{4pt}

If $A = \begin{pmatrix} 1 & 2 \\ 2 & -1 \end{pmatrix}$, then which ONE of the following is $A^8$?

\par\vspace{6pt}\noindent
\begin{tabular}{@{}p{0.98\linewidth}@{}}
$\square$\ \textbf{(A)} $\begin{pmatrix} 25 & 0 \\ 0 & 25 \end{pmatrix}$ \\
\addlinespace[3pt]
$\square$\ \textbf{(B)} $\begin{pmatrix} 125 & 0 \\ 0 & 125 \end{pmatrix}$ \\
\addlinespace[3pt]
$\square$\ \textbf{(C)} $\begin{pmatrix} 625 & 0 \\ 0 & 625 \end{pmatrix}$ \\
\addlinespace[3pt]
$\square$\ \textbf{(D)} $\begin{pmatrix} 3125 & 0 \\ 0 & 3125 \end{pmatrix}$
\end{tabular}

\tcblower

\textbf{Official Key:} \textbf{\large C} \hfill \textit{\small [Medium]}
\vspace{2pt}\hrule height 0.3pt\vspace{3pt}
\textbf{Topic:} Linear Algebra \\
\textbf{Subtopic:} \textit{Matrices and matrix powers}
\vspace{2pt}\hrule height 0.3pt\vspace{3pt}
\textbf{Guide / Hint:} \\
{\small Compute $A^2 = \begin{pmatrix} 1 & 2 \\ 2 & -1 \end{pmatrix}^2 = \begin{pmatrix} 5 & 0 \\ 0 & 5 \end{pmatrix} = 5I$. Then $A^8 = (A^2)^4 = (5I)^4 = 5^4 I = 625I = \begin{pmatrix} 625 & 0 \\ 0 & 625 \end{pmatrix}$.}

\end{tcolorbox}
\vspace{6pt}\hrule height 0.35pt\vspace{6pt}
\noindent
\begin{tcolorbox}[
    blanker,
    sidebyside,
    sidebyside align=top,
    lefthand width=0.55\textwidth,
    top=4pt, bottom=4pt, left=0pt, right=0pt,
    middle=16pt,
    underlay={
        \draw[dashed, black!80, line width=0.8pt, dash pattern=on 4pt off 3pt]
            ([xshift=0.55\textwidth+8pt]frame.north west) --
            ([xshift=0.55\textwidth+8pt]frame.south west);
    }
]
\textbf{Q.12} \quad {\footnotesize\textsc{[MCQ $\cdot$ 1 Mark(s) $\cdot$ Neg: 0.33]}} \hfill {\tiny\texttt{2025\_CS2\_CS\_Q12}}
\vspace{2pt}\hrule height 0.3pt\vspace{4pt}

The value of $x$ such that $x > 1$, satisfying the equation $\int_1^x t \ln t \, dt = \frac{1}{4}$ is

\par\vspace{6pt}\noindent
\begin{tabular}{@{}p{0.48\linewidth}@{\hspace{0.04\linewidth}}p{0.48\linewidth}@{}}
$\square$\ \textbf{(A)} $\sqrt{e}$ & $\square$\ \textbf{(B)} $e$ \\
\addlinespace[3pt]
$\square$\ \textbf{(C)} $e^2$ & $\square$\ \textbf{(D)} $e^{-1}$
\end{tabular}

\tcblower

\textbf{Official Key:} \textbf{\large A} \hfill \textit{\small [Medium]}
\vspace{2pt}\hrule height 0.3pt\vspace{3pt}
\textbf{Topic:} Calculus \\
\textbf{Subtopic:} \textit{Integration and definite integrals}
\vspace{2pt}\hrule height 0.3pt\vspace{3pt}
\textbf{Guide / Hint:} \\
{\small Using integration by parts, $\int_1^x t \ln t \, dt = \left[\frac{t^2}{2} \ln t - \frac{t^2}{4}\right]_1^x = \frac{x^2}{2} \ln x - \frac{x^2}{4} + \frac{1}{4}$. Equating this to $\frac{1}{4}$ gives $\frac{x^2}{4}(2\ln x - 1) = 0$. Since $x > 1$, we have $2\ln x = 1 \implies x = e^{1/2} = \sqrt{e}$.}

\end{tcolorbox}
\vspace{6pt}\hrule height 0.35pt\vspace{6pt}
\noindent
\begin{tcolorbox}[
    blanker,
    sidebyside,
    sidebyside align=top,
    lefthand width=0.55\textwidth,
    top=4pt, bottom=4pt, left=0pt, right=0pt,
    middle=16pt,
    underlay={
        \draw[dashed, black!80, line width=0.8pt, dash pattern=on 4pt off 3pt]
            ([xshift=0.55\textwidth+8pt]frame.north west) --
            ([xshift=0.55\textwidth+8pt]frame.south west);
    }
]
\textbf{Q.13} \quad {\footnotesize\textsc{[MCQ $\cdot$ 1 Mark(s) $\cdot$ Neg: 0.33]}} \hfill {\tiny\texttt{2025\_CS2\_CS\_Q13}}
\vspace{2pt}\hrule height 0.3pt\vspace{4pt}

Consider a binary tree $T$ in which every node has either zero or two children. Let $n > 0$ be the number of nodes in $T$. Which ONE of the following is the number of nodes in $T$ that have exactly two children?

\par\vspace{6pt}\noindent
\begin{tabular}{@{}p{0.48\linewidth}@{\hspace{0.04\linewidth}}p{0.48\linewidth}@{}}
$\square$\ \textbf{(A)} $\frac{n-2}{2}$ & $\square$\ \textbf{(B)} $\frac{n-1}{2}$ \\
\addlinespace[3pt]
$\square$\ \textbf{(C)} $\frac{n}{2}$ & $\square$\ \textbf{(D)} $\frac{n+1}{2}$
\end{tabular}

\tcblower

\textbf{Official Key:} \textbf{\large B} \hfill \textit{\small [Easy]}
\vspace{2pt}\hrule height 0.3pt\vspace{3pt}
\textbf{Topic:} Programming and Data Structures \\
\textbf{Subtopic:} \textit{Trees and strictly binary trees}
\vspace{2pt}\hrule height 0.3pt\vspace{3pt}
\textbf{Guide / Hint:} \\
{\small In any proper (strictly) binary tree with $L$ leaves and $I$ internal nodes (nodes with 2 children), $L = I + 1$. Since $n = L + I = 2I + 1$, the number of internal nodes is $I = \frac{n-1}{2}$.}

\end{tcolorbox}
\vspace{6pt}\hrule height 0.35pt\vspace{6pt}
\noindent
\begin{tcolorbox}[
    blanker,
    sidebyside,
    sidebyside align=top,
    lefthand width=0.55\textwidth,
    top=4pt, bottom=4pt, left=0pt, right=0pt,
    middle=16pt,
    underlay={
        \draw[dashed, black!80, line width=0.8pt, dash pattern=on 4pt off 3pt]
            ([xshift=0.55\textwidth+8pt]frame.north west) --
            ([xshift=0.55\textwidth+8pt]frame.south west);
    }
]
\textbf{Q.14} \quad {\footnotesize\textsc{[MCQ $\cdot$ 1 Mark(s) $\cdot$ Neg: 0.33]}} \hfill {\tiny\texttt{2025\_CS2\_CS\_Q14}}
\vspace{2pt}\hrule height 0.3pt\vspace{4pt}

Let $L$, $M$, and $N$ be non-singular matrices of order 3 satisfying the equations $L^2 = L^{-1}$, $M = L^8$ and $N = L^2$. Which ONE of the following is the value of the determinant of $(M - N)$?

\par\vspace{6pt}\noindent
\begin{tabular}{@{}p{0.48\linewidth}@{\hspace{0.04\linewidth}}p{0.48\linewidth}@{}}
$\square$\ \textbf{(A)} 0 & $\square$\ \textbf{(B)} 1 \\
\addlinespace[3pt]
$\square$\ \textbf{(C)} 2 & $\square$\ \textbf{(D)} 3
\end{tabular}

\tcblower

\textbf{Official Key:} \textbf{\large A} \hfill \textit{\small [Medium]}
\vspace{2pt}\hrule height 0.3pt\vspace{3pt}
\textbf{Topic:} Linear Algebra \\
\textbf{Subtopic:} \textit{Matrices and determinants}
\vspace{2pt}\hrule height 0.3pt\vspace{3pt}
\textbf{Guide / Hint:} \\
{\small From $L^2 = L^{-1}$, multiplying by $L$ gives $L^3 = I$. Then $M = L^8 = (L^3)^2 L^2 = L^2$, and $N = L^2$. Hence $M - N = L^2 - L^2 = O$, whose determinant is $\det(O) = 0$.}

\end{tcolorbox}
\vspace{6pt}\hrule height 0.35pt\vspace{6pt}
\noindent
\begin{tcolorbox}[
    blanker,
    sidebyside,
    sidebyside align=top,
    lefthand width=0.55\textwidth,
    top=4pt, bottom=4pt, left=0pt, right=0pt,
    middle=16pt,
    underlay={
        \draw[dashed, black!80, line width=0.8pt, dash pattern=on 4pt off 3pt]
            ([xshift=0.55\textwidth+8pt]frame.north west) --
            ([xshift=0.55\textwidth+8pt]frame.south west);
    }
]
\textbf{Q.15} \quad {\footnotesize\textsc{[MCQ $\cdot$ 1 Mark(s) $\cdot$ Neg: 0.33]}} \hfill {\tiny\texttt{2025\_CS2\_CS\_Q15}}
\vspace{2pt}\hrule height 0.3pt\vspace{4pt}

Let $P(x)$ be an arbitrary predicate over the domain of natural numbers. Which ONE of the following statements is TRUE?

\par\vspace{6pt}\noindent
\begin{tabular}{@{}p{0.98\linewidth}@{}}
$\square$\ \textbf{(A)} $(P(0) \land (\forall x [P(x) \implies P(x+1)])) \implies (\forall x P(x))$ \\
\addlinespace[3pt]
$\square$\ \textbf{(B)} $(P(0) \land (\forall x [P(x) \implies P(x-1)])) \implies (\forall x P(x))$ \\
\addlinespace[3pt]
$\square$\ \textbf{(C)} $(P(1000) \land (\forall x [P(x) \implies P(x-1)])) \implies (\forall x P(x))$ \\
\addlinespace[3pt]
$\square$\ \textbf{(D)} $(P(1000) \land (\forall x [P(x) \implies P(x+1)])) \implies (\forall x P(x))$
\end{tabular}

\tcblower

\textbf{Official Key:} \textbf{\large A} \hfill \textit{\small [Easy]}
\vspace{2pt}\hrule height 0.3pt\vspace{3pt}
\textbf{Topic:} Discrete Mathematics \\
\textbf{Subtopic:} \textit{Propositional and first order logic: mathematical induction}
\vspace{2pt}\hrule height 0.3pt\vspace{3pt}
\textbf{Guide / Hint:} \\
{\small Option A is the formal statement of the Principle of Mathematical Induction over $\mathbb{N}$: if the base case $P(0)$ holds and for all $x$, $P(x) \implies P(x+1)$, then $\forall x, P(x)$ is true.}

\end{tcolorbox}
\vspace{6pt}\hrule height 0.35pt\vspace{6pt}
\noindent
\begin{tcolorbox}[
    blanker,
    sidebyside,
    sidebyside align=top,
    lefthand width=0.55\textwidth,
    top=4pt, bottom=4pt, left=0pt, right=0pt,
    middle=16pt,
    underlay={
        \draw[dashed, black!80, line width=0.8pt, dash pattern=on 4pt off 3pt]
            ([xshift=0.55\textwidth+8pt]frame.north west) --
            ([xshift=0.55\textwidth+8pt]frame.south west);
    }
]
\textbf{Q.16} \quad {\footnotesize\textsc{[MCQ $\cdot$ 1 Mark(s) $\cdot$ Neg: 0.33]}} \hfill {\tiny\texttt{2025\_CS2\_CS\_Q16}}
\vspace{2pt}\hrule height 0.3pt\vspace{4pt}

Consider the following statements:
(i) Address Resolution Protocol (ARP) provides a mapping from an IP address to the corresponding hardware (link-layer) address.
(ii) A single TCP segment from a sender S to a receiver R cannot carry both data from S to R and acknowledgement for a segment from R to S.
Which ONE of the following is CORRECT?

\par\vspace{6pt}\noindent
\begin{tabular}{@{}p{0.48\linewidth}@{\hspace{0.04\linewidth}}p{0.48\linewidth}@{}}
$\square$\ \textbf{(A)} Both (i) and (ii) are TRUE & $\square$\ \textbf{(B)} (i) is TRUE and (ii) is FALSE \\
\addlinespace[3pt]
$\square$\ \textbf{(C)} (i) is FALSE and (ii) is TRUE & $\square$\ \textbf{(D)} Both (i) and (ii) are FALSE
\end{tabular}

\tcblower

\textbf{Official Key:} \textbf{\large B} \hfill \textit{\small [Easy]}
\vspace{2pt}\hrule height 0.3pt\vspace{3pt}
\textbf{Topic:} Computer Networks \\
\textbf{Subtopic:} \textit{Basics of IP support protocols: ARP and TCP headers}
\vspace{2pt}\hrule height 0.3pt\vspace{3pt}
\textbf{Guide / Hint:} \\
{\small ARP resolves an IPv4 address to a physical MAC (link-layer) address, so (i) is true. A TCP segment supports piggybacking, where data and acknowledgement flags/numbers are sent in the same segment, so (ii) is false.}

\end{tcolorbox}
\vspace{6pt}\hrule height 0.35pt\vspace{6pt}
\noindent
\begin{tcolorbox}[
    blanker,
    sidebyside,
    sidebyside align=top,
    lefthand width=0.55\textwidth,
    top=4pt, bottom=4pt, left=0pt, right=0pt,
    middle=16pt,
    underlay={
        \draw[dashed, black!80, line width=0.8pt, dash pattern=on 4pt off 3pt]
            ([xshift=0.55\textwidth+8pt]frame.north west) --
            ([xshift=0.55\textwidth+8pt]frame.south west);
    }
]
\textbf{Q.17} \quad {\footnotesize\textsc{[MCQ $\cdot$ 1 Mark(s) $\cdot$ Neg: 0.33]}} \hfill {\tiny\texttt{2025\_CS2\_CS\_Q17}}
\vspace{2pt}\hrule height 0.3pt\vspace{4pt}

Consider the routing protocols given in List I and the names given in List II:
\par\medskip
\begin{center}\small
\begin{tabular}{ll}
\toprule
List I & List II \\
\midrule
(i) Distance vector routing & (a) Bellman-Ford \\
(ii) Link state routing & (b) Dijkstra \\
\bottomrule
\end{tabular}\end{center}
\par\medskip
For matching of items in List I with those in List II, which ONE of the following options is CORRECT?

\par\vspace{6pt}\noindent
\begin{tabular}{@{}p{0.48\linewidth}@{\hspace{0.04\linewidth}}p{0.48\linewidth}@{}}
$\square$\ \textbf{(A)} (i) – (a) and (ii) – (b) & $\square$\ \textbf{(B)} (i) – (a) and (ii) – (a) \\
\addlinespace[3pt]
$\square$\ \textbf{(C)} (i) – (b) and (ii) – (a) & $\square$\ \textbf{(D)} (i) – (b) and (ii) – (b)
\end{tabular}

\tcblower

\textbf{Official Key:} \textbf{\large A} \hfill \textit{\small [Easy]}
\vspace{2pt}\hrule height 0.3pt\vspace{3pt}
\textbf{Topic:} Computer Networks \\
\textbf{Subtopic:} \textit{Routing protocols: distance vector and link state}
\vspace{2pt}\hrule height 0.3pt\vspace{3pt}
\textbf{Guide / Hint:} \\
{\small Distance Vector routing is based on the distributed Bellman-Ford algorithm, whereas Link State routing uses Dijkstra's shortest path algorithm.}

\end{tcolorbox}
\vspace{6pt}\hrule height 0.35pt\vspace{6pt}
\noindent
\begin{tcolorbox}[
    blanker,
    sidebyside,
    sidebyside align=top,
    lefthand width=0.55\textwidth,
    top=4pt, bottom=4pt, left=0pt, right=0pt,
    middle=16pt,
    underlay={
        \draw[dashed, black!80, line width=0.8pt, dash pattern=on 4pt off 3pt]
            ([xshift=0.55\textwidth+8pt]frame.north west) --
            ([xshift=0.55\textwidth+8pt]frame.south west);
    }
]
\textbf{Q.18} \quad {\footnotesize\textsc{[MCQ $\cdot$ 1 Mark(s) $\cdot$ Neg: 0.33]}} \hfill {\tiny\texttt{2025\_CS2\_CS\_Q18}}
\vspace{2pt}\hrule height 0.3pt\vspace{4pt}

A machine receives an IPv4 datagram. The protocol field of the IPv4 header has the protocol number of a protocol X. Which ONE of the following is NOT a possible candidate for X?

\par\vspace{6pt}\noindent
\begin{tabular}{@{}p{0.98\linewidth}@{}}
$\square$\ \textbf{(A)} Internet Control Message Protocol (ICMP) \\
\addlinespace[3pt]
$\square$\ \textbf{(B)} Internet Group Management Protocol (IGMP) \\
\addlinespace[3pt]
$\square$\ \textbf{(C)} Open Shortest Path First (OSPF) \\
\addlinespace[3pt]
$\square$\ \textbf{(D)} Routing Information Protocol (RIP)
\end{tabular}

\tcblower

\textbf{Official Key:} \textbf{\large D} \hfill \textit{\small [Medium]}
\vspace{2pt}\hrule height 0.3pt\vspace{3pt}
\textbf{Topic:} Computer Networks \\
\textbf{Subtopic:} \textit{Concept of layering: OSI and TCP/IP protocol stacks}
\vspace{2pt}\hrule height 0.3pt\vspace{3pt}
\textbf{Guide / Hint:} \\
{\small The IPv4 Protocol field identifies the payload carried directly in the IP datagram: ICMP (protocol 1), IGMP (protocol 2), and OSPF (protocol 89). In contrast, RIP runs as an application-layer process encapsulated inside UDP (port 520).}

\end{tcolorbox}
\vspace{6pt}\hrule height 0.35pt\vspace{6pt}
\noindent
\begin{tcolorbox}[
    blanker,
    sidebyside,
    sidebyside align=top,
    lefthand width=0.55\textwidth,
    top=4pt, bottom=4pt, left=0pt, right=0pt,
    middle=16pt,
    underlay={
        \draw[dashed, black!80, line width=0.8pt, dash pattern=on 4pt off 3pt]
            ([xshift=0.55\textwidth+8pt]frame.north west) --
            ([xshift=0.55\textwidth+8pt]frame.south west);
    }
]
\textbf{Q.19} \quad {\footnotesize\textsc{[MCQ $\cdot$ 1 Mark(s) $\cdot$ Neg: 0.33]}} \hfill {\tiny\texttt{2025\_CS2\_CS\_Q19}}
\vspace{2pt}\hrule height 0.3pt\vspace{4pt}

Consider the following C program:
\begin{lstlisting}[language=C, basicstyle=\footnotesize\ttfamily, frame=single, rulecolor=\color{black!30}, tabsize=2, breaklines=true]
#include <stdio.h>
void stringcopy(char *, char *);
int main() {
    char a[30] = "@#Hello World!";
    stringcopy(a, a + 2);
    printf("%s\n", a);
    return 0;
}
void stringcopy(char *s, char *t) {
    while(*t)
        *s++ = *t++;
}
\end{lstlisting}
Which ONE of the following will be the output of the program?

\par\vspace{6pt}\noindent
\begin{tabular}{@{}p{0.48\linewidth}@{\hspace{0.04\linewidth}}p{0.48\linewidth}@{}}
$\square$\ \textbf{(A)} @\#Hello World! & $\square$\ \textbf{(B)} Hello World! \\
\addlinespace[3pt]
$\square$\ \textbf{(C)} ello World! & $\square$\ \textbf{(D)} Hello World!d!
\end{tabular}

\tcblower

\textbf{Official Key:} \textbf{\large D} \hfill \textit{\small [Medium]}
\vspace{2pt}\hrule height 0.3pt\vspace{3pt}
\textbf{Topic:} Programming and Data Structures \\
\textbf{Subtopic:} \textit{Programming in C: pointers and strings}
\vspace{2pt}\hrule height 0.3pt\vspace{3pt}
\textbf{Guide / Hint:} \\
{\small `stringcopy(a, a + 2)` copies string `"Hello World!"` into the start of array `a`, overwriting the first 12 characters. The copy stops when `*t` hits null terminator without copying the null terminator over `d!`, leaving trailing `d!\0`, resulting in `"Hello World!d!"`.}

\end{tcolorbox}
\vspace{6pt}\hrule height 0.35pt\vspace{6pt}
\noindent
\begin{tcolorbox}[
    blanker,
    sidebyside,
    sidebyside align=top,
    lefthand width=0.55\textwidth,
    top=4pt, bottom=4pt, left=0pt, right=0pt,
    middle=16pt,
    underlay={
        \draw[dashed, black!80, line width=0.8pt, dash pattern=on 4pt off 3pt]
            ([xshift=0.55\textwidth+8pt]frame.north west) --
            ([xshift=0.55\textwidth+8pt]frame.south west);
    }
]
\textbf{Q.20} \quad {\footnotesize\textsc{[MCQ $\cdot$ 1 Mark(s) $\cdot$ Neg: 0.33]}} \hfill {\tiny\texttt{2025\_CS2\_CS\_Q20}}
\vspace{2pt}\hrule height 0.3pt\vspace{4pt}

Consider an unordered list of $N$ distinct integers. What is the minimum number of element comparisons required to find an integer in the list that is NOT the largest in the list?

\par\vspace{6pt}\noindent
\begin{tabular}{@{}p{0.48\linewidth}@{\hspace{0.04\linewidth}}p{0.48\linewidth}@{}}
$\square$\ \textbf{(A)} 1 & $\square$\ \textbf{(B)} $N - 1$ \\
\addlinespace[3pt]
$\square$\ \textbf{(C)} $N$ & $\square$\ \textbf{(D)} $2N - 1$
\end{tabular}

\tcblower

\textbf{Official Key:} \textbf{\large A} \hfill \textit{\small [Easy]}
\vspace{2pt}\hrule height 0.3pt\vspace{3pt}
\textbf{Topic:} Algorithms \\
\textbf{Subtopic:} \textit{Searching and comparison complexity}
\vspace{2pt}\hrule height 0.3pt\vspace{3pt}
\textbf{Guide / Hint:} \\
{\small To find any element that is not the maximum, compare any two distinct elements in 1 comparison; the smaller of the two is guaranteed not to be the maximum.}

\end{tcolorbox}
\vspace{6pt}\hrule height 0.35pt\vspace{6pt}
\noindent
\begin{tcolorbox}[
    blanker,
    sidebyside,
    sidebyside align=top,
    lefthand width=0.55\textwidth,
    top=4pt, bottom=4pt, left=0pt, right=0pt,
    middle=16pt,
    underlay={
        \draw[dashed, black!80, line width=0.8pt, dash pattern=on 4pt off 3pt]
            ([xshift=0.55\textwidth+8pt]frame.north west) --
            ([xshift=0.55\textwidth+8pt]frame.south west);
    }
]
\textbf{Q.21} \quad {\footnotesize\textsc{[MCQ $\cdot$ 1 Mark(s) $\cdot$ Neg: 0.33]}} \hfill {\tiny\texttt{2025\_CS2\_CS\_Q21}}
\vspace{2pt}\hrule height 0.3pt\vspace{4pt}

Consider the following statements about the use of backpatching in a compiler for intermediate code generation:
(I) Backpatching can be used to generate code for Boolean expression in one pass.
(II) Backpatching can be used to generate code for flow-of-control statements in one pass.
Which ONE of the following options is CORRECT?

\par\vspace{6pt}\noindent
\begin{tabular}{@{}p{0.98\linewidth}@{}}
$\square$\ \textbf{(A)} Only (I) is correct. \\
\addlinespace[3pt]
$\square$\ \textbf{(B)} Only (II) is correct. \\
\addlinespace[3pt]
$\square$\ \textbf{(C)} Both (I) and (II) are correct. \\
\addlinespace[3pt]
$\square$\ \textbf{(D)} Neither (I) nor (II) is correct.
\end{tabular}

\tcblower

\textbf{Official Key:} \textbf{\large C} \hfill \textit{\small [Easy]}
\vspace{2pt}\hrule height 0.3pt\vspace{3pt}
\textbf{Topic:} Compiler Design \\
\textbf{Subtopic:} \textit{Intermediate code generation and backpatching}
\vspace{2pt}\hrule height 0.3pt\vspace{3pt}
\textbf{Guide / Hint:} \\
{\small Backpatching is a standard technique used in single-pass compilers to synthesize target labels for forward jumps in both Boolean expressions and flow-of-control statements.}

\end{tcolorbox}
\vspace{6pt}\hrule height 0.35pt\vspace{6pt}
\noindent
\begin{tcolorbox}[
    blanker,
    sidebyside,
    sidebyside align=top,
    lefthand width=0.55\textwidth,
    top=4pt, bottom=4pt, left=0pt, right=0pt,
    middle=16pt,
    underlay={
        \draw[dashed, black!80, line width=0.8pt, dash pattern=on 4pt off 3pt]
            ([xshift=0.55\textwidth+8pt]frame.north west) --
            ([xshift=0.55\textwidth+8pt]frame.south west);
    }
]
\textbf{Q.22} \quad {\footnotesize\textsc{[MCQ $\cdot$ 1 Mark(s) $\cdot$ Neg: 0.33]}} \hfill {\tiny\texttt{2025\_CS2\_CS\_Q22}}
\vspace{2pt}\hrule height 0.3pt\vspace{4pt}

Given the following syntax directed translation rules:
- Rule 1: $R \to AB \;\{ B.i = R.i - 1; \; A.i = B.i; \; R.i = A.i + 1; \}$
- Rule 2: $P \to CD \;\{ P.i = C.i + D.i; \; D.i = C.i + 2; \}$
- Rule 3: $Q \to EF \;\{ Q.i = E.i + F.i; \}$
\par\medskip
Which ONE is the CORRECT option among the following?

\par\vspace{6pt}\noindent
\begin{tabular}{@{}p{0.98\linewidth}@{}}
$\square$\ \textbf{(A)} Rule 1 is S-attributed and L-attributed; Rule 2 is S-attributed and not L-attributed; Rule 3 is neither S-attributed nor L-attributed \\
\addlinespace[3pt]
$\square$\ \textbf{(B)} Rule 1 is neither S-attributed nor L-attributed; Rule 2 is S-attributed and L-attributed; Rule 3 is S-attributed and L-attributed \\
\addlinespace[3pt]
$\square$\ \textbf{(C)} Rule 1 is neither S-attributed nor L-attributed; Rule 2 is not S-attributed and is L-attributed; Rule 3 is S-attributed and L-attributed \\
\addlinespace[3pt]
$\square$\ \textbf{(D)} Rule 1 is S-attributed and not L-attributed; Rule 2 is not S-attributed and is L-attributed; Rule 3 is S-attributed and L-attributed
\end{tabular}

\tcblower

\textbf{Official Key:} \textbf{\large C} \hfill \textit{\small [Hard]}
\vspace{2pt}\hrule height 0.3pt\vspace{3pt}
\textbf{Topic:} Compiler Design \\
\textbf{Subtopic:} \textit{Syntax-directed translation: S-attributed and L-attributed definitions}
\vspace{2pt}\hrule height 0.3pt\vspace{3pt}
\textbf{Guide / Hint:} \\
{\small In Rule 1, $R.i$ is computed from children (synthesized-like) but defined circularly with inherited attributes, making it neither. In Rule 2, attributes depend only on left siblings, satisfying L-attributed definition. In Rule 3, $Q.i$ depends only on children, which is S-attributed (and hence L-attributed).}

\end{tcolorbox}
\vspace{6pt}\hrule height 0.35pt\vspace{6pt}
\noindent
\begin{tcolorbox}[
    blanker,
    sidebyside,
    sidebyside align=top,
    lefthand width=0.55\textwidth,
    top=4pt, bottom=4pt, left=0pt, right=0pt,
    middle=16pt,
    underlay={
        \draw[dashed, black!80, line width=0.8pt, dash pattern=on 4pt off 3pt]
            ([xshift=0.55\textwidth+8pt]frame.north west) --
            ([xshift=0.55\textwidth+8pt]frame.south west);
    }
]
\textbf{Q.23} \quad {\footnotesize\textsc{[MCQ $\cdot$ 1 Mark(s) $\cdot$ Neg: 0.33]}} \hfill {\tiny\texttt{2025\_CS2\_CS\_Q23}}
\vspace{2pt}\hrule height 0.3pt\vspace{4pt}

Consider a network that uses Ethernet and IPv4. Assume that IPv4 headers do not use any options field. Each Ethernet frame can carry a maximum of 1500 bytes in its data field. A UDP segment is transmitted. The payload (data) in the UDP segment is 7488 bytes.
Which ONE of the following choices has the CORRECT total number of fragments transmitted and the size of the last fragment including IPv4 header?

\par\vspace{6pt}\noindent
\begin{tabular}{@{}p{0.48\linewidth}@{\hspace{0.04\linewidth}}p{0.48\linewidth}@{}}
$\square$\ \textbf{(A)} 5 fragments, 1488 bytes & $\square$\ \textbf{(B)} 6 fragments, 88 bytes \\
\addlinespace[3pt]
$\square$\ \textbf{(C)} 6 fragments, 108 bytes & $\square$\ \textbf{(D)} 6 fragments, 116 bytes
\end{tabular}

\tcblower

\textbf{Official Key:} \textbf{\large D} \hfill \textit{\small [Medium]}
\vspace{2pt}\hrule height 0.3pt\vspace{3pt}
\textbf{Topic:} Computer Networks \\
\textbf{Subtopic:} \textit{Fragmentation and IP addressing}
\vspace{2pt}\hrule height 0.3pt\vspace{3pt}
\textbf{Guide / Hint:} \\
{\small UDP payload is 7488 bytes + 8 bytes UDP header = 7496 bytes IP payload. Max IP payload per Ethernet fragment is $\lfloor(1500 - 20)/8\rfloor \times 8 = 1480$ bytes. Fragment sizes are 5 fragments of 1480 bytes ($5 \times 1480 = 7400$), leaving $7496 - 7400 = 96$ bytes for the 6th fragment, so total size is $96 + 20 = 116$ bytes.}

\end{tcolorbox}
\vspace{6pt}\hrule height 0.35pt\vspace{6pt}
\noindent
\begin{tcolorbox}[
    blanker,
    sidebyside,
    sidebyside align=top,
    lefthand width=0.55\textwidth,
    top=4pt, bottom=4pt, left=0pt, right=0pt,
    middle=16pt,
    underlay={
        \draw[dashed, black!80, line width=0.8pt, dash pattern=on 4pt off 3pt]
            ([xshift=0.55\textwidth+8pt]frame.north west) --
            ([xshift=0.55\textwidth+8pt]frame.south west);
    }
]
\textbf{Q.24} \quad {\footnotesize\textsc{[MCQ $\cdot$ 1 Mark(s) $\cdot$ Neg: 0.33]}} \hfill {\tiny\texttt{2025\_CS2\_CS\_Q24}}
\vspace{2pt}\hrule height 0.3pt\vspace{4pt}

Which ONE of the following languages is accepted by a deterministic pushdown automaton?

\par\vspace{6pt}\noindent
\begin{tabular}{@{}p{0.98\linewidth}@{}}
$\square$\ \textbf{(A)} Any regular language. \\
\addlinespace[3pt]
$\square$\ \textbf{(B)} Any context-free language. \\
\addlinespace[3pt]
$\square$\ \textbf{(C)} Any language accepted by a non-deterministic pushdown automaton. \\
\addlinespace[3pt]
$\square$\ \textbf{(D)} Any decidable language.
\end{tabular}

\tcblower

\textbf{Official Key:} \textbf{\large A} \hfill \textit{\small [Easy]}
\vspace{2pt}\hrule height 0.3pt\vspace{3pt}
\textbf{Topic:} Theory of Computation \\
\textbf{Subtopic:} \textit{Context-free grammars and push-down automata}
\vspace{2pt}\hrule height 0.3pt\vspace{3pt}
\textbf{Guide / Hint:} \\
{\small Every regular language is recognized by a Deterministic Finite Automaton (DFA), which can be directly simulated by a DPDA without using stack operations.}

\end{tcolorbox}
\vspace{6pt}\hrule height 0.35pt\vspace{6pt}
\noindent
\begin{tcolorbox}[
    blanker,
    sidebyside,
    sidebyside align=top,
    lefthand width=0.55\textwidth,
    top=4pt, bottom=4pt, left=0pt, right=0pt,
    middle=16pt,
    underlay={
        \draw[dashed, black!80, line width=0.8pt, dash pattern=on 4pt off 3pt]
            ([xshift=0.55\textwidth+8pt]frame.north west) --
            ([xshift=0.55\textwidth+8pt]frame.south west);
    }
]
\textbf{Q.25} \quad {\footnotesize\textsc{[MCQ $\cdot$ 1 Mark(s) $\cdot$ Neg: 0.33]}} \hfill {\tiny\texttt{2025\_CS2\_CS\_Q25}}
\vspace{2pt}\hrule height 0.3pt\vspace{4pt}

Let $G_1, G_2$ be Context Free Grammars (CFGs) and $R$ be a regular expression. For a grammar $G$, let $L(G)$ denote the language generated by $G$. Which ONE among the following questions is decidable?

\par\vspace{6pt}\noindent
\begin{tabular}{@{}p{0.98\linewidth}@{}}
$\square$\ \textbf{(A)} Is $L(G_1) = L(G_2)$? \\
\addlinespace[3pt]
$\square$\ \textbf{(B)} Is $L(G_1) \cap L(G_2) = \emptyset$? \\
\addlinespace[3pt]
$\square$\ \textbf{(C)} Is $L(G_1) = L(R)$? \\
\addlinespace[3pt]
$\square$\ \textbf{(D)} Is $L(G_1) = \emptyset$?
\end{tabular}

\tcblower

\textbf{Official Key:} \textbf{\large D} \hfill \textit{\small [Easy]}
\vspace{2pt}\hrule height 0.3pt\vspace{3pt}
\textbf{Topic:} Theory of Computation \\
\textbf{Subtopic:} \textit{Turing machines and undecidability}
\vspace{2pt}\hrule height 0.3pt\vspace{3pt}
\textbf{Guide / Hint:} \\
{\small For Context-Free Grammars, emptiness ($L(G) = \emptyset$) is decidable via finding generating non-terminals, whereas equivalence ($L(G_1) = L(G_2)$) and disjointness ($L(G_1) \cap L(G_2) = \emptyset$) are undecidable.}

\end{tcolorbox}
\vspace{6pt}\hrule height 0.35pt\vspace{6pt}
\noindent
\begin{tcolorbox}[
    blanker,
    sidebyside,
    sidebyside align=top,
    lefthand width=0.55\textwidth,
    top=4pt, bottom=4pt, left=0pt, right=0pt,
    middle=16pt,
    underlay={
        \draw[dashed, black!80, line width=0.8pt, dash pattern=on 4pt off 3pt]
            ([xshift=0.55\textwidth+8pt]frame.north west) --
            ([xshift=0.55\textwidth+8pt]frame.south west);
    }
]
\textbf{Q.26} \quad {\footnotesize\textsc{[MCQ $\cdot$ 1 Mark(s) $\cdot$ Neg: 0.33]}} \hfill {\tiny\texttt{2025\_CS2\_CS\_Q26}}
\vspace{2pt}\hrule height 0.3pt\vspace{4pt}

Processes $P_1, P_2, P_3, P_4$ arrive in that order at times 0, 1, 2, and 8 milliseconds respectively, and have execution times of 10, 13, 6, and 9 milliseconds respectively. Shortest Remaining Time First (SRTF) algorithm is used as the CPU scheduling policy. Ignore context switching times. Which ONE of the following correctly gives the average turnaround time of the four processes in milliseconds?

\par\vspace{6pt}\noindent
\begin{tabular}{@{}p{0.48\linewidth}@{\hspace{0.04\linewidth}}p{0.48\linewidth}@{}}
$\square$\ \textbf{(A)} 22 & $\square$\ \textbf{(B)} 15 \\
\addlinespace[3pt]
$\square$\ \textbf{(C)} 37 & $\square$\ \textbf{(D)} 19
\end{tabular}

\tcblower

\textbf{Official Key:} \textbf{\large D} \hfill \textit{\small [Medium]}
\vspace{2pt}\hrule height 0.3pt\vspace{3pt}
\textbf{Topic:} Operating System \\
\textbf{Subtopic:} \textit{CPU and I/O scheduling: SRTF}
\vspace{2pt}\hrule height 0.3pt\vspace{3pt}
\textbf{Guide / Hint:} \\
{\small Using preemptive SRTF: $P_1$ runs [0,2], $P_3$ runs [2,8] (finishes at 8, TAT = $8-2=6$). Then $P_4$ (remaining 9) runs [8,17] (finishes at 17, TAT = $17-8=9$). $P_1$ (remaining 8) runs [17,25] (finishes at 25, TAT = $25-0=25$). $P_2$ runs [25,38] (finishes at 38, TAT = $38-1=37$). Average TAT = $\frac{25 + 37 + 6 + 9}{4} = \frac{76}{4} = 19$ ms.}

\end{tcolorbox}
\vspace{6pt}\hrule height 0.35pt\vspace{6pt}
\noindent
\begin{tcolorbox}[
    blanker,
    sidebyside,
    sidebyside align=top,
    lefthand width=0.55\textwidth,
    top=4pt, bottom=4pt, left=0pt, right=0pt,
    middle=16pt,
    underlay={
        \draw[dashed, black!80, line width=0.8pt, dash pattern=on 4pt off 3pt]
            ([xshift=0.55\textwidth+8pt]frame.north west) --
            ([xshift=0.55\textwidth+8pt]frame.south west);
    }
]
\textbf{Q.27} \quad {\footnotesize\textsc{[MSQ $\cdot$ 1 Mark(s) $\cdot$ No Neg]}} \hfill {\tiny\texttt{2025\_CS2\_CS\_Q27}}
\vspace{2pt}\hrule height 0.3pt\vspace{4pt}

An audit of a banking transactions system has found that on an earlier occasion, two joint holders of account $A$ attempted simultaneous transfers of Rs. 10000 each from account $A$ to account $B$. Both transactions read the same value, Rs. 11000, as the initial balance in $A$ and were allowed to go through. $B$ was credited Rs. 10000 twice. $A$ was debited only once and ended up with a balance of Rs. 1000. Which of the following properties is/are certain to have been violated by the system?

\par\vspace{6pt}\noindent
\begin{tabular}{@{}p{0.48\linewidth}@{\hspace{0.04\linewidth}}p{0.48\linewidth}@{}}
$\square$\ \textbf{(A)} Atomicity & $\square$\ \textbf{(B)} Consistency \\
\addlinespace[3pt]
$\square$\ \textbf{(C)} Isolation & $\square$\ \textbf{(D)} Durability
\end{tabular}

\tcblower

\textbf{Official Key:} \textbf{\large B;C} \hfill \textit{\small [Medium]}
\vspace{2pt}\hrule height 0.3pt\vspace{3pt}
\textbf{Topic:} Databases \\
\textbf{Subtopic:} \textit{Transactions and concurrency control: ACID properties}
\vspace{2pt}\hrule height 0.3pt\vspace{3pt}
\textbf{Guide / Hint:} \\
{\small Both transactions read the same uncommitted state and overwrote each other (lost update anomaly), violating Isolation. The final balance violates the database invariant (starting 11000 debited $2 \times 10000$ cannot yield 1000), violating Consistency.}

\end{tcolorbox}
\vspace{6pt}\hrule height 0.35pt\vspace{6pt}
\noindent
\begin{tcolorbox}[
    blanker,
    sidebyside,
    sidebyside align=top,
    lefthand width=0.55\textwidth,
    top=4pt, bottom=4pt, left=0pt, right=0pt,
    middle=16pt,
    underlay={
        \draw[dashed, black!80, line width=0.8pt, dash pattern=on 4pt off 3pt]
            ([xshift=0.55\textwidth+8pt]frame.north west) --
            ([xshift=0.55\textwidth+8pt]frame.south west);
    }
]
\textbf{Q.28} \quad {\footnotesize\textsc{[MSQ $\cdot$ 1 Mark(s) $\cdot$ No Neg]}} \hfill {\tiny\texttt{2025\_CS2\_CS\_Q28}}
\vspace{2pt}\hrule height 0.3pt\vspace{4pt}

Which of the following is/are part of an Instruction Set Architecture of a processor?

\par\vspace{6pt}\noindent
\begin{tabular}{@{}p{0.98\linewidth}@{}}
$\square$\ \textbf{(A)} The size of the cache memory \\
\addlinespace[3pt]
$\square$\ \textbf{(B)} The clock frequency of the processor \\
\addlinespace[3pt]
$\square$\ \textbf{(C)} The number of cache memory levels \\
\addlinespace[3pt]
$\square$\ \textbf{(D)} The total number of registers
\end{tabular}

\tcblower

\textbf{Official Key:} \textbf{\large D} \hfill \textit{\small [Easy]}
\vspace{2pt}\hrule height 0.3pt\vspace{3pt}
\textbf{Topic:} Computer Organization and Architecture \\
\textbf{Subtopic:} \textit{Machine instructions and addressing modes: ISA}
\vspace{2pt}\hrule height 0.3pt\vspace{3pt}
\textbf{Guide / Hint:} \\
{\small The Instruction Set Architecture (ISA) defines programmer-visible specifications, including the register set/count and instruction encoding. Cache sizes, frequency, and microarchitecture levels are implementation details, not part of the ISA.}

\end{tcolorbox}
\vspace{6pt}\hrule height 0.35pt\vspace{6pt}
\noindent
\begin{tcolorbox}[
    blanker,
    sidebyside,
    sidebyside align=top,
    lefthand width=0.55\textwidth,
    top=4pt, bottom=4pt, left=0pt, right=0pt,
    middle=16pt,
    underlay={
        \draw[dashed, black!80, line width=0.8pt, dash pattern=on 4pt off 3pt]
            ([xshift=0.55\textwidth+8pt]frame.north west) --
            ([xshift=0.55\textwidth+8pt]frame.south west);
    }
]
\textbf{Q.29} \quad {\footnotesize\textsc{[MSQ $\cdot$ 1 Mark(s) $\cdot$ No Neg]}} \hfill {\tiny\texttt{2025\_CS2\_CS\_Q29}}
\vspace{2pt}\hrule height 0.3pt\vspace{4pt}

Which of the following statements regarding Breadth First Search (BFS) and Depth First Search (DFS) on an undirected simple graph $G$ is/are TRUE?

\par\vspace{6pt}\noindent
\begin{tabular}{@{}p{0.98\linewidth}@{}}
$\square$\ \textbf{(A)} A DFS tree of $G$ is a Shortest Path tree of $G$. \\
\addlinespace[3pt]
$\square$\ \textbf{(B)} Every non-tree edge of $G$ with respect to a DFS tree is a forward/back edge. \\
\addlinespace[3pt]
$\square$\ \textbf{(C)} If $(u, v)$ is a non-tree edge of $G$ with respect to a BFS tree, then the distances from the source vertex $s$ to $u$ and $v$ in the BFS tree are within $\pm 1$ of each other. \\
\addlinespace[3pt]
$\square$\ \textbf{(D)} Both BFS and DFS can be used to find the connected components of $G$.
\end{tabular}

\tcblower

\textbf{Official Key:} \textbf{\large B;C;D} \hfill \textit{\small [Medium]}
\vspace{2pt}\hrule height 0.3pt\vspace{3pt}
\textbf{Topic:} Algorithms \\
\textbf{Subtopic:} \textit{Graph traversals: BFS and DFS}
\vspace{2pt}\hrule height 0.3pt\vspace{3pt}
\textbf{Guide / Hint:} \\
{\small In an undirected graph, non-tree DFS edges can only be back edges (B is true). In a BFS tree, cross edges can only exist between the same or adjacent levels, so distances differ by at most $\pm 1$ (C is true). Both can compute connected components (D is true).}

\end{tcolorbox}
\vspace{6pt}\hrule height 0.35pt\vspace{6pt}
\noindent
\begin{tcolorbox}[
    blanker,
    sidebyside,
    sidebyside align=top,
    lefthand width=0.55\textwidth,
    top=4pt, bottom=4pt, left=0pt, right=0pt,
    middle=16pt,
    underlay={
        \draw[dashed, black!80, line width=0.8pt, dash pattern=on 4pt off 3pt]
            ([xshift=0.55\textwidth+8pt]frame.north west) --
            ([xshift=0.55\textwidth+8pt]frame.south west);
    }
]
\textbf{Q.30} \quad {\footnotesize\textsc{[MSQ $\cdot$ 1 Mark(s) $\cdot$ No Neg]}} \hfill {\tiny\texttt{2025\_CS2\_CS\_Q30}}
\vspace{2pt}\hrule height 0.3pt\vspace{4pt}

Consider the two lists List I and List II given below:
\par\medskip
\begin{center}\small
\begin{tabular}{ll}
\toprule
List I & List II \\
\midrule
(i) Context free languages & (a) Closed under union \\
(ii) Recursive languages & (b) Not closed under complementation \\
(iii) Regular languages & (c) Closed under intersection \\
\bottomrule
\end{tabular}\end{center}
\par\medskip
For matching of items in List I with those in List II, which of the following option(s) is/are CORRECT?

\par\vspace{6pt}\noindent
\begin{tabular}{@{}p{0.98\linewidth}@{}}
$\square$\ \textbf{(A)} (i) – (a), (ii) – (b), and (iii) – (c) \\
\addlinespace[3pt]
$\square$\ \textbf{(B)} (i) – (b), (ii) – (a), and (iii) – (c) \\
\addlinespace[3pt]
$\square$\ \textbf{(C)} (i) – (b), (ii) – (c), and (iii) – (a) \\
\addlinespace[3pt]
$\square$\ \textbf{(D)} (i) – (a), (ii) – (c), and (iii) – (b)
\end{tabular}

\tcblower

\textbf{Official Key:} \textbf{\large B;C} \hfill \textit{\small [Medium]}
\vspace{2pt}\hrule height 0.3pt\vspace{3pt}
\textbf{Topic:} Theory of Computation \\
\textbf{Subtopic:} \textit{Regular and context-free languages: closure properties}
\vspace{2pt}\hrule height 0.3pt\vspace{3pt}
\textbf{Guide / Hint:} \\
{\small Context-free languages are not closed under complementation (i - b) or intersection, but are closed under union. Recursive languages are closed under union and complementation and intersection (ii - a or c). Regular languages are closed under all three.}

\end{tcolorbox}
\vspace{6pt}\hrule height 0.35pt\vspace{6pt}
\noindent
\begin{tcolorbox}[
    blanker,
    sidebyside,
    sidebyside align=top,
    lefthand width=0.55\textwidth,
    top=4pt, bottom=4pt, left=0pt, right=0pt,
    middle=16pt,
    underlay={
        \draw[dashed, black!80, line width=0.8pt, dash pattern=on 4pt off 3pt]
            ([xshift=0.55\textwidth+8pt]frame.north west) --
            ([xshift=0.55\textwidth+8pt]frame.south west);
    }
]
\textbf{Q.31} \quad {\footnotesize\textsc{[MSQ $\cdot$ 1 Mark(s) $\cdot$ No Neg]}} \hfill {\tiny\texttt{2025\_CS2\_CS\_Q31}}
\vspace{2pt}\hrule height 0.3pt\vspace{4pt}

Consider the logic circuit diagram for function $F$. Which is/are the CORRECT option(s) for the output function $F$?

\par\vspace{6pt}\noindent
\begin{tabular}{@{}p{0.98\linewidth}@{}}
$\square$\ \textbf{(A)} $\overline{X \overline{Y}}$ \\
\addlinespace[3pt]
$\square$\ \textbf{(B)} $\overline{X} + \overline{Y} + X\overline{Y}$ \\
\addlinespace[3pt]
$\square$\ \textbf{(C)} $\overline{X \overline{Y}} + \overline{X} + X\overline{Y}$ \\
\addlinespace[3pt]
$\square$\ \textbf{(D)} $X + \overline{Y}$
\end{tabular}

\tcblower

\textbf{Official Key:} \textbf{\large A;B;C} \hfill \textit{\small [Medium]}
\vspace{2pt}\hrule height 0.3pt\vspace{3pt}
\textbf{Topic:} Digital Logic \\
\textbf{Subtopic:} \textit{Combinational and sequential circuits: logic gates}
\vspace{2pt}\hrule height 0.3pt\vspace{3pt}
\textbf{Guide / Hint:} \\
{\small From Boolean reduction, $\overline{X \cdot \overline{Y}} = \overline{X} + Y$. Applying consensus and absorption theorems yields equivalent representations including $\overline{X} + \overline{Y} + X\overline{Y}$ and $\overline{X \cdot \overline{Y}} + \overline{X} + X\overline{Y}$.}

\end{tcolorbox}
\vspace{6pt}\hrule height 0.35pt\vspace{6pt}
\noindent
\begin{tcolorbox}[
    blanker,
    sidebyside,
    sidebyside align=top,
    lefthand width=0.55\textwidth,
    top=4pt, bottom=4pt, left=0pt, right=0pt,
    middle=16pt,
    underlay={
        \draw[dashed, black!80, line width=0.8pt, dash pattern=on 4pt off 3pt]
            ([xshift=0.55\textwidth+8pt]frame.north west) --
            ([xshift=0.55\textwidth+8pt]frame.south west);
    }
]
\textbf{Q.32} \quad {\footnotesize\textsc{[NAT $\cdot$ 1 Mark(s) $\cdot$ No Neg]}} \hfill {\tiny\texttt{2025\_CS2\_CS\_Q32}}
\vspace{2pt}\hrule height 0.3pt\vspace{4pt}

The following two signed 2’s complement numbers (multiplicand $M$ and multiplier $Q$) are being multiplied using Booth’s algorithm:
$M$: 1100 1101 1110 1101 and $Q$: 1010 0100 1010 1010
The total number of addition and subtraction operations to be performed is \_\_\_\_\_\_\_\_\_\_\_. (Answer in integer)

\par\vspace{6pt}\noindent\textbf{Your Answer:} \fbox{\rule{0pt}{1.3ex}\hspace{3cm}}

\tcblower

\textbf{Official Key:} \textbf{\large 13 to 13} \hfill \textit{\small [Hard]}
\vspace{2pt}\hrule height 0.3pt\vspace{3pt}
\textbf{Topic:} Computer Organization and Architecture \\
\textbf{Subtopic:} \textit{Number representations and computer arithmetic: Booth's algorithm}
\vspace{2pt}\hrule height 0.3pt\vspace{3pt}
\textbf{Guide / Hint:} \\
{\small In Booth's algorithm, additions and subtractions occur whenever there is a transition between adjacent bits $(q_i, q_{i-1})$ of multiplier $Q$ with $q_{-1}=0$. Scanning $Q = 1010010010101010$ with $q_{-1}=0$ produces exactly 13 bit-transitions ($01$ or $10$).}

\end{tcolorbox}
\vspace{6pt}\hrule height 0.35pt\vspace{6pt}
\noindent
\begin{tcolorbox}[
    blanker,
    sidebyside,
    sidebyside align=top,
    lefthand width=0.55\textwidth,
    top=4pt, bottom=4pt, left=0pt, right=0pt,
    middle=16pt,
    underlay={
        \draw[dashed, black!80, line width=0.8pt, dash pattern=on 4pt off 3pt]
            ([xshift=0.55\textwidth+8pt]frame.north west) --
            ([xshift=0.55\textwidth+8pt]frame.south west);
    }
]
\textbf{Q.33} \quad {\footnotesize\textsc{[NAT $\cdot$ 1 Mark(s) $\cdot$ No Neg]}} \hfill {\tiny\texttt{2025\_CS2\_CS\_Q33}}
\vspace{2pt}\hrule height 0.3pt\vspace{4pt}

Consider the following C code segment:
\begin{lstlisting}[language=C, basicstyle=\footnotesize\ttfamily, frame=single, rulecolor=\color{black!30}, tabsize=2, breaklines=true]
int x=126, y=105;
do {
    if(x > y) x = x - y;
    else y = y - x;
} while(x != y);
printf("%d", x);
\end{lstlisting}
The output of the given C code segment is \_\_\_\_\_\_\_\_. (Answer in integer)

\par\vspace{6pt}\noindent\textbf{Your Answer:} \fbox{\rule{0pt}{1.3ex}\hspace{3cm}}

\tcblower

\textbf{Official Key:} \textbf{\large 21 to 21} \hfill \textit{\small [Easy]}
\vspace{2pt}\hrule height 0.3pt\vspace{3pt}
\textbf{Topic:} Programming and Data Structures \\
\textbf{Subtopic:} \textit{Programming in C: loops and Euclidean GCD}
\vspace{2pt}\hrule height 0.3pt\vspace{3pt}
\textbf{Guide / Hint:} \\
{\small The do-while loop implements the subtraction-based Euclidean algorithm for greatest common divisor: $\gcd(x, y) = \gcd(126, 105) = 21$.}

\end{tcolorbox}
\vspace{6pt}\hrule height 0.35pt\vspace{6pt}
\noindent
\begin{tcolorbox}[
    blanker,
    sidebyside,
    sidebyside align=top,
    lefthand width=0.55\textwidth,
    top=4pt, bottom=4pt, left=0pt, right=0pt,
    middle=16pt,
    underlay={
        \draw[dashed, black!80, line width=0.8pt, dash pattern=on 4pt off 3pt]
            ([xshift=0.55\textwidth+8pt]frame.north west) --
            ([xshift=0.55\textwidth+8pt]frame.south west);
    }
]
\textbf{Q.34} \quad {\footnotesize\textsc{[NAT $\cdot$ 1 Mark(s) $\cdot$ No Neg]}} \hfill {\tiny\texttt{2025\_CS2\_CS\_Q34}}
\vspace{2pt}\hrule height 0.3pt\vspace{4pt}

In a 4-bit ripple counter, if the period of the waveform at the last flip-flop is $64\,\mu\text{s}$, then the frequency of the ripple counter in kHz is \_\_\_\_\_\_\_\_. (Answer in integer)

\par\vspace{6pt}\noindent\textbf{Your Answer:} \fbox{\rule{0pt}{1.3ex}\hspace{3cm}}

\tcblower

\textbf{Official Key:} \textbf{\large 250 to 250} \hfill \textit{\small [Medium]}
\vspace{2pt}\hrule height 0.3pt\vspace{3pt}
\textbf{Topic:} Digital Logic \\
\textbf{Subtopic:} \textit{Combinational and sequential circuits: ripple counters}
\vspace{2pt}\hrule height 0.3pt\vspace{3pt}
\textbf{Guide / Hint:} \\
{\small A 4-bit ripple counter divides the clock frequency by $2^4 = 16$. The clock period is $T_{clk} = \frac{T_{out}}{16} = \frac{64\,\mu\text{s}}{16} = 4\,\mu\text{s}$, giving clock frequency $f = \frac{1}{4\,\mu\text{s}} = 250\,\text{kHz}$.}

\end{tcolorbox}
\vspace{6pt}\hrule height 0.35pt\vspace{6pt}
\noindent
\begin{tcolorbox}[
    blanker,
    sidebyside,
    sidebyside align=top,
    lefthand width=0.55\textwidth,
    top=4pt, bottom=4pt, left=0pt, right=0pt,
    middle=16pt,
    underlay={
        \draw[dashed, black!80, line width=0.8pt, dash pattern=on 4pt off 3pt]
            ([xshift=0.55\textwidth+8pt]frame.north west) --
            ([xshift=0.55\textwidth+8pt]frame.south west);
    }
]
\textbf{Q.35} \quad {\footnotesize\textsc{[NAT $\cdot$ 1 Mark(s) $\cdot$ No Neg]}} \hfill {\tiny\texttt{2025\_CS2\_CS\_Q35}}
\vspace{2pt}\hrule height 0.3pt\vspace{4pt}

Suppose the values 10, $-4$, 15, 30, 20, 5, 60, 19 are inserted in that order into an initially empty binary search tree. Let $T$ be the resulting binary search tree. The number of edges in the path from the node containing 19 to the root node of $T$ is \_\_\_\_\_\_\_\_\_\_\_. (Answer in integer)

\par\vspace{6pt}\noindent\textbf{Your Answer:} \fbox{\rule{0pt}{1.3ex}\hspace{3cm}}

\tcblower

\textbf{Official Key:} \textbf{\large 4 to 4} \hfill \textit{\small [Easy]}
\vspace{2pt}\hrule height 0.3pt\vspace{3pt}
\textbf{Topic:} Programming and Data Structures \\
\textbf{Subtopic:} \textit{Trees: binary search trees}
\vspace{2pt}\hrule height 0.3pt\vspace{3pt}
\textbf{Guide / Hint:} \\
{\small Insert keys into BST: 10 is root $\to$ 15 goes right $\to$ 30 goes right $\to$ 20 goes left $\to$ 19 goes left of 20. The path from root is $10 \to 15 \to 30 \to 20 \to 19$, which contains 4 edges.}

\end{tcolorbox}
\vspace{6pt}\hrule height 0.35pt\vspace{6pt}
\noindent
\begin{tcolorbox}[
    blanker,
    sidebyside,
    sidebyside align=top,
    lefthand width=0.55\textwidth,
    top=4pt, bottom=4pt, left=0pt, right=0pt,
    middle=16pt,
    underlay={
        \draw[dashed, black!80, line width=0.8pt, dash pattern=on 4pt off 3pt]
            ([xshift=0.55\textwidth+8pt]frame.north west) --
            ([xshift=0.55\textwidth+8pt]frame.south west);
    }
]
\textbf{Q.36} \quad {\footnotesize\textsc{[MCQ $\cdot$ 2 Mark(s) $\cdot$ Neg: 0.66]}} \hfill {\tiny\texttt{2025\_CS2\_CS\_Q36}}
\vspace{2pt}\hrule height 0.3pt\vspace{4pt}

Suppose we are transmitting frames between two nodes using Stop-and-Wait protocol. The frame size is 3000 bits. Which ONE of the following gives the efficiency of the channel?

\par\vspace{6pt}\noindent
\begin{tabular}{@{}p{0.48\linewidth}@{\hspace{0.04\linewidth}}p{0.48\linewidth}@{}}
$\square$\ \textbf{(A)} 88.23 & $\square$\ \textbf{(B)} 93.75 \\
\addlinespace[3pt]
$\square$\ \textbf{(C)} 85.44 & $\square$\ \textbf{(D)} 66.67
\end{tabular}

\tcblower

\textbf{Official Key:} \textbf{\large A} \hfill \textit{\small [Medium]}
\vspace{2pt}\hrule height 0.3pt\vspace{3pt}
\textbf{Topic:} Computer Networks \\
\textbf{Subtopic:} \textit{Data link layer: flow control and Stop-and-Wait efficiency}
\vspace{2pt}\hrule height 0.3pt\vspace{3pt}
\textbf{Guide / Hint:} \\
{\small Efficiency in Stop-and-Wait is $\eta = \frac{T_t}{T_t + 2 T_p}$. Substituting given transmission and propagation delays yields $\eta = 88.23\%$.}

\end{tcolorbox}
\vspace{6pt}\hrule height 0.35pt\vspace{6pt}
\noindent
\begin{tcolorbox}[
    blanker,
    sidebyside,
    sidebyside align=top,
    lefthand width=0.55\textwidth,
    top=4pt, bottom=4pt, left=0pt, right=0pt,
    middle=16pt,
    underlay={
        \draw[dashed, black!80, line width=0.8pt, dash pattern=on 4pt off 3pt]
            ([xshift=0.55\textwidth+8pt]frame.north west) --
            ([xshift=0.55\textwidth+8pt]frame.south west);
    }
]
\textbf{Q.37} \quad {\footnotesize\textsc{[MCQ $\cdot$ 2 Mark(s) $\cdot$ Neg: 0.66]}} \hfill {\tiny\texttt{2025\_CS2\_CS\_Q37}}
\vspace{2pt}\hrule height 0.3pt\vspace{4pt}

Let $G$ be an edge-weighted undirected graph with positive edge weights. Suppose a positive constant $\alpha$ is added to the weight of every edge in $G$. Which ONE of the following statements is TRUE?

\par\vspace{6pt}\noindent
\begin{tabular}{@{}p{0.98\linewidth}@{}}
$\square$\ \textbf{(A)} Every MST remains an MST, and every SP remains an SP. \\
\addlinespace[3pt]
$\square$\ \textbf{(B)} MSTs need not remain MSTs, and every SP remains an SP. \\
\addlinespace[3pt]
$\square$\ \textbf{(C)} Every MST remains an MST, and SPs need not remain SPs. \\
\addlinespace[3pt]
$\square$\ \textbf{(D)} MSTs need not remain MSTs, and SPs need not remain SPs.
\end{tabular}

\tcblower

\textbf{Official Key:} \textbf{\large C} \hfill \textit{\small [Medium]}
\vspace{2pt}\hrule height 0.3pt\vspace{3pt}
\textbf{Topic:} Algorithms \\
\textbf{Subtopic:} \textit{Graph algorithms: minimum spanning trees and shortest paths}
\vspace{2pt}\hrule height 0.3pt\vspace{3pt}
\textbf{Guide / Hint:} \\
{\small Adding a constant $\alpha$ to all edge weights preserves the relative order of all spanning trees because every spanning tree has exactly $V - 1$ edges (so all tree costs increase by $(V-1)\alpha$). However, shortest paths with varying numbers of edges can change because paths with more edges receive a larger cumulative penalty.}

\end{tcolorbox}
\vspace{6pt}\hrule height 0.35pt\vspace{6pt}
\noindent
\begin{tcolorbox}[
    blanker,
    sidebyside,
    sidebyside align=top,
    lefthand width=0.55\textwidth,
    top=4pt, bottom=4pt, left=0pt, right=0pt,
    middle=16pt,
    underlay={
        \draw[dashed, black!80, line width=0.8pt, dash pattern=on 4pt off 3pt]
            ([xshift=0.55\textwidth+8pt]frame.north west) --
            ([xshift=0.55\textwidth+8pt]frame.south west);
    }
]
\textbf{Q.38} \quad {\footnotesize\textsc{[MCQ $\cdot$ 2 Mark(s) $\cdot$ Neg: 0.66]}} \hfill {\tiny\texttt{2025\_CS2\_CS\_Q38}}
\vspace{2pt}\hrule height 0.3pt\vspace{4pt}

A meld operation on two instances of a data structure combines them into one single instance of the same data structure. What are the best worst-case time complexities of meld on: $P$: circular doubly linked lists, $Q$: standard binary heaps, and $R$: arrays?

\par\vspace{6pt}\noindent
\begin{tabular}{@{}p{0.98\linewidth}@{}}
$\square$\ \textbf{(A)} P: $\Theta(1)$, Q: $\Theta(n)$, R: $\Theta(n)$ \\
\addlinespace[3pt]
$\square$\ \textbf{(B)} P: $\Theta(1)$, Q: $\Theta(n\log n)$, R: $\Theta(n)$ \\
\addlinespace[3pt]
$\square$\ \textbf{(C)} P: $\Theta(n)$, Q: $\Theta(n\log n)$, R: $\Theta(n^2)$ \\
\addlinespace[3pt]
$\square$\ \textbf{(D)} P: $\Theta(1)$, Q: $\Theta(n)$, R: $\Theta(n\log n)$
\end{tabular}

\tcblower

\textbf{Official Key:} \textbf{\large A} \hfill \textit{\small [Hard]}
\vspace{2pt}\hrule height 0.3pt\vspace{3pt}
\textbf{Topic:} Programming and Data Structures \\
\textbf{Subtopic:} \textit{Data structures: heaps, linked lists and merge operations}
\vspace{2pt}\hrule height 0.3pt\vspace{3pt}
\textbf{Guide / Hint:} \\
{\small Merging two circular linked lists takes $\Theta(1)$ by swapping pointers. Merging two binary heaps takes $\Theta(n)$ using Floyd's build-heap algorithm. Hence, time complexities are $P: \Theta(1), Q: \Theta(n), R: \Theta(n)$.}

\end{tcolorbox}
\vspace{6pt}\hrule height 0.35pt\vspace{6pt}
\noindent
\begin{tcolorbox}[
    blanker,
    sidebyside,
    sidebyside align=top,
    lefthand width=0.55\textwidth,
    top=4pt, bottom=4pt, left=0pt, right=0pt,
    middle=16pt,
    underlay={
        \draw[dashed, black!80, line width=0.8pt, dash pattern=on 4pt off 3pt]
            ([xshift=0.55\textwidth+8pt]frame.north west) --
            ([xshift=0.55\textwidth+8pt]frame.south west);
    }
]
\textbf{Q.39} \quad {\footnotesize\textsc{[MCQ $\cdot$ 2 Mark(s) $\cdot$ Neg: 0.66]}} \hfill {\tiny\texttt{2025\_CS2\_CS\_Q39}}
\vspace{2pt}\hrule height 0.3pt\vspace{4pt}

For a direct-mapped cache, 4 bits are used for the tag field and 12 bits are used to index into a cache block. What are the cache data size and tag directory size respectively?

\par\vspace{6pt}\noindent
\begin{tabular}{@{}p{0.48\linewidth}@{\hspace{0.04\linewidth}}p{0.48\linewidth}@{}}
$\square$\ \textbf{(A)} 64 KB and 4 KB & $\square$\ \textbf{(B)} 128 KB and 16 KB \\
\addlinespace[3pt]
$\square$\ \textbf{(C)} 64 KB and 8 KB & $\square$\ \textbf{(D)} 128 KB and 6 KB
\end{tabular}

\tcblower

\textbf{Official Key:} \textbf{\large A} \hfill \textit{\small [Medium]}
\vspace{2pt}\hrule height 0.3pt\vspace{3pt}
\textbf{Topic:} Computer Organization and Architecture \\
\textbf{Subtopic:} \textit{Memory hierarchy: cache memory mapping}
\vspace{2pt}\hrule height 0.3pt\vspace{3pt}
\textbf{Guide / Hint:} \\
{\small Number of cache lines is $2^{12} = 4096$. With 4 tag bits and physical address specification, cache block size is determined, yielding total cache size of 64 KB and tag directory size of 4 KB.}

\end{tcolorbox}
\vspace{6pt}\hrule height 0.35pt\vspace{6pt}
\noindent
\begin{tcolorbox}[
    blanker,
    sidebyside,
    sidebyside align=top,
    lefthand width=0.55\textwidth,
    top=4pt, bottom=4pt, left=0pt, right=0pt,
    middle=16pt,
    underlay={
        \draw[dashed, black!80, line width=0.8pt, dash pattern=on 4pt off 3pt]
            ([xshift=0.55\textwidth+8pt]frame.north west) --
            ([xshift=0.55\textwidth+8pt]frame.south west);
    }
]
\textbf{Q.40} \quad {\footnotesize\textsc{[MCQ $\cdot$ 2 Mark(s) $\cdot$ Neg: 0.66]}} \hfill {\tiny\texttt{2025\_CS2\_CS\_Q40}}
\vspace{2pt}\hrule height 0.3pt\vspace{4pt}

Given a Context-Free Grammar $G$:
$$\begin{aligned} S &\to Aa \mid bAc \mid dc \mid bda \\ A &\to d \end{aligned}$$
Which ONE of the following statements is CORRECT?

\par\vspace{6pt}\noindent
\begin{tabular}{@{}p{0.98\linewidth}@{}}
$\square$\ \textbf{(A)} $G$ is neither LALR(1) nor SLR(1) \\
\addlinespace[3pt]
$\square$\ \textbf{(B)} $G$ is CLR(1), not LALR(1) \\
\addlinespace[3pt]
$\square$\ \textbf{(C)} $G$ is LALR(1), not SLR(1) \\
\addlinespace[3pt]
$\square$\ \textbf{(D)} $G$ is LALR(1), also SLR(1)
\end{tabular}

\tcblower

\textbf{Official Key:} \textbf{\large C} \hfill \textit{\small [Hard]}
\vspace{2pt}\hrule height 0.3pt\vspace{3pt}
\textbf{Topic:} Compiler Design \\
\textbf{Subtopic:} \textit{Parsing: LR, LALR, and SLR parsers}
\vspace{2pt}\hrule height 0.3pt\vspace{3pt}
\textbf{Guide / Hint:} \\
{\small In SLR(1), FOLLOW sets create reduce-reduce conflicts on lookahead $c$. In LALR(1), the merged states retain distinct LR(1) lookaheads that avoid conflicts, so the grammar is LALR(1) but not SLR(1).}

\end{tcolorbox}
\vspace{6pt}\hrule height 0.35pt\vspace{6pt}
\noindent
\begin{tcolorbox}[
    blanker,
    sidebyside,
    sidebyside align=top,
    lefthand width=0.55\textwidth,
    top=4pt, bottom=4pt, left=0pt, right=0pt,
    middle=16pt,
    underlay={
        \draw[dashed, black!80, line width=0.8pt, dash pattern=on 4pt off 3pt]
            ([xshift=0.55\textwidth+8pt]frame.north west) --
            ([xshift=0.55\textwidth+8pt]frame.south west);
    }
]
\textbf{Q.41} \quad {\footnotesize\textsc{[MCQ $\cdot$ 2 Mark(s) $\cdot$ Neg: 0.66]}} \hfill {\tiny\texttt{2025\_CS2\_CS\_Q41}}
\vspace{2pt}\hrule height 0.3pt\vspace{4pt}

An array $A$ of length $n$ with distinct elements is said to be bitonic if there is an index $1 \le i \le n$ such that $A[1..i]$ is strictly increasing and $A[i..n]$ is strictly decreasing. What is the optimal time complexity to search for a key in a bitonic array?

\par\vspace{6pt}\noindent
\begin{tabular}{@{}p{0.48\linewidth}@{\hspace{0.04\linewidth}}p{0.48\linewidth}@{}}
$\square$\ \textbf{(A)} $\Theta(n)$ & $\square$\ \textbf{(B)} $\Theta(1)$ \\
\addlinespace[3pt]
$\square$\ \textbf{(C)} $\Theta(\log^2 n)$ & $\square$\ \textbf{(D)} $\Theta(\log n)$
\end{tabular}

\tcblower

\textbf{Official Key:} \textbf{\large D} \hfill \textit{\small [Medium]}
\vspace{2pt}\hrule height 0.3pt\vspace{3pt}
\textbf{Topic:} Algorithms \\
\textbf{Subtopic:} \textit{Searching and divide-and-conquer: bitonic array search}
\vspace{2pt}\hrule height 0.3pt\vspace{3pt}
\textbf{Guide / Hint:} \\
{\small In a bitonic array, the maximum element can be found in $\Theta(\log n)$ time using binary search on mid-elements. Once split into increasing and decreasing halves, binary search on each half takes $\Theta(\log n)$, giving overall time complexity $\Theta(\log n)$.}

\end{tcolorbox}
\vspace{6pt}\hrule height 0.35pt\vspace{6pt}
\noindent
\begin{tcolorbox}[
    blanker,
    sidebyside,
    sidebyside align=top,
    lefthand width=0.55\textwidth,
    top=4pt, bottom=4pt, left=0pt, right=0pt,
    middle=16pt,
    underlay={
        \draw[dashed, black!80, line width=0.8pt, dash pattern=on 4pt off 3pt]
            ([xshift=0.55\textwidth+8pt]frame.north west) --
            ([xshift=0.55\textwidth+8pt]frame.south west);
    }
]
\textbf{Q.42} \quad {\footnotesize\textsc{[MSQ $\cdot$ 2 Mark(s) $\cdot$ No Neg]}} \hfill {\tiny\texttt{2025\_CS2\_CS\_Q42}}
\vspace{2pt}\hrule height 0.3pt\vspace{4pt}

Let $\mathcal{F}$ be the set of all functions from $\{1, \dots, n\}$ to $\{0, 1\}$. Define the binary relation $\preceq$ on $\mathcal{F}$ as follows: $\forall f, g \in \mathcal{F}, f \preceq g \iff \forall i \in \{1, \dots, n\}, f(i) \le g(i)$. Which of the following is/are TRUE?

\par\vspace{6pt}\noindent
\begin{tabular}{@{}p{0.98\linewidth}@{}}
$\square$\ \textbf{(A)} $\preceq$ is a symmetric relation \\
\addlinespace[3pt]
$\square$\ \textbf{(B)} $(\mathcal{F}, \preceq)$ is a partial order \\
\addlinespace[3pt]
$\square$\ \textbf{(C)} $(\mathcal{F}, \preceq)$ is a lattice \\
\addlinespace[3pt]
$\square$\ \textbf{(D)} $\preceq$ is an equivalence relation
\end{tabular}

\tcblower

\textbf{Official Key:} \textbf{\large B;C} \hfill \textit{\small [Hard]}
\vspace{2pt}\hrule height 0.3pt\vspace{3pt}
\textbf{Topic:} Discrete Mathematics \\
\textbf{Subtopic:} \textit{Sets, relations, functions, partial orders and lattices}
\vspace{2pt}\hrule height 0.3pt\vspace{3pt}
\textbf{Guide / Hint:} \\
{\small The pointwise ordering $f \preceq g \iff \forall i, f(i) \le g(i)$ is reflexive, antisymmetric, and transitive, forming a partial order (poset). Furthermore, every pair of functions has a unique supremum $\max(f, g)$ and infimum $\min(f, g)$, so $(\mathcal{F}, \preceq)$ is a distributive lattice.}

\end{tcolorbox}
\vspace{6pt}\hrule height 0.35pt\vspace{6pt}
\noindent
\begin{tcolorbox}[
    blanker,
    sidebyside,
    sidebyside align=top,
    lefthand width=0.55\textwidth,
    top=4pt, bottom=4pt, left=0pt, right=0pt,
    middle=16pt,
    underlay={
        \draw[dashed, black!80, line width=0.8pt, dash pattern=on 4pt off 3pt]
            ([xshift=0.55\textwidth+8pt]frame.north west) --
            ([xshift=0.55\textwidth+8pt]frame.south west);
    }
]
\textbf{Q.43} \quad {\footnotesize\textsc{[MSQ $\cdot$ 2 Mark(s) $\cdot$ No Neg]}} \hfill {\tiny\texttt{2025\_CS2\_CS\_Q43}}
\vspace{2pt}\hrule height 0.3pt\vspace{4pt}

Given the Karnaugh Map for a Boolean function $F(w, x, y, z)$, which one or more of the following Boolean expressions represent $F$?

\par\vspace{6pt}\noindent
\begin{tabular}{@{}p{0.98\linewidth}@{}}
$\square$\ \textbf{(A)} $\overline{w}\overline{x}\overline{y}\overline{z} + w\overline{x}\overline{y}\overline{z} + \overline{w}\overline{x}y\overline{z} + w\overline{x}y\overline{z} + xz$ \\
\addlinespace[3pt]
$\square$\ \textbf{(B)} $\overline{w}\overline{x}\overline{y}\overline{z} + \overline{w}\overline{x}y\overline{z} + w\overline{x}yz + xz$ \\
\addlinespace[3pt]
$\square$\ \textbf{(C)} $\overline{w}\overline{x}\overline{y}\overline{z} + w\overline{x}\overline{y}\overline{z} + w\overline{x}\overline{y}z + xz$ \\
\addlinespace[3pt]
$\square$\ \textbf{(D)} $\overline{x}\overline{z} + xz$
\end{tabular}

\tcblower

\textbf{Official Key:} \textbf{\large A;D} \hfill \textit{\small [Medium]}
\vspace{2pt}\hrule height 0.3pt\vspace{3pt}
\textbf{Topic:} Digital Logic \\
\textbf{Subtopic:} \textit{Boolean algebra and minimization: Karnaugh Map}
\vspace{2pt}\hrule height 0.3pt\vspace{3pt}
\textbf{Guide / Hint:} \\
{\small Grouping the 1s in the K-map: the four corners and the central groupings give prime implicants $\overline{x}\overline{z} + xz$. Both the expanded minterm expression and minimal SOP are valid representations.}

\end{tcolorbox}
\vspace{6pt}\hrule height 0.35pt\vspace{6pt}
\noindent
\begin{tcolorbox}[
    blanker,
    sidebyside,
    sidebyside align=top,
    lefthand width=0.55\textwidth,
    top=4pt, bottom=4pt, left=0pt, right=0pt,
    middle=16pt,
    underlay={
        \draw[dashed, black!80, line width=0.8pt, dash pattern=on 4pt off 3pt]
            ([xshift=0.55\textwidth+8pt]frame.north west) --
            ([xshift=0.55\textwidth+8pt]frame.south west);
    }
]
\textbf{Q.44} \quad {\footnotesize\textsc{[MSQ $\cdot$ 2 Mark(s) $\cdot$ No Neg]}} \hfill {\tiny\texttt{2025\_CS2\_CS\_Q44}}
\vspace{2pt}\hrule height 0.3pt\vspace{4pt}

Consider a system of linear equations $PX = Q$ where $P \in \mathbb{R}^{3 \times 3}$ and $Q \in \mathbb{R}^{3 \times 1}$. Suppose $P$ has an LU decomposition $P = LU$, where $L$ is unit lower triangular and $U$ is upper triangular. Which of the following statements is/are TRUE?

\par\vspace{6pt}\noindent
\begin{tabular}{@{}p{0.98\linewidth}@{}}
$\square$\ \textbf{(A)} The system $PX = Q$ can be solved by first solving $LY = Q$ and then $UX = Y$. \\
\addlinespace[3pt]
$\square$\ \textbf{(B)} If $P$ is invertible, then both $L$ and $U$ are invertible. \\
\addlinespace[3pt]
$\square$\ \textbf{(C)} If $P$ is singular, then at least one of the diagonal elements of $U$ is zero. \\
\addlinespace[3pt]
$\square$\ \textbf{(D)} If $P$ is symmetric, then both $L$ and $U$ are symmetric.
\end{tabular}

\tcblower

\textbf{Official Key:} \textbf{\large A;B;C} \hfill \textit{\small [Hard]}
\vspace{2pt}\hrule height 0.3pt\vspace{3pt}
\textbf{Topic:} Linear Algebra \\
\textbf{Subtopic:} \textit{System of linear equations and LU decomposition}
\vspace{2pt}\hrule height 0.3pt\vspace{3pt}
\textbf{Guide / Hint:} \\
{\small For $PX = Q$ with $P = LU$, solving $LY = Q$ by forward substitution and then $UX = Y$ by backward substitution solves the system. Since $\det(P) = \det(L)\det(U)$ and $L$ has 1s on the diagonal ($\\det(L)=1$), $P$ is invertible iff $U$ is invertible, and singular iff a diagonal entry of $U$ is 0.}

\end{tcolorbox}
\vspace{6pt}\hrule height 0.35pt\vspace{6pt}
\noindent
\begin{tcolorbox}[
    blanker,
    sidebyside,
    sidebyside align=top,
    lefthand width=0.55\textwidth,
    top=4pt, bottom=4pt, left=0pt, right=0pt,
    middle=16pt,
    underlay={
        \draw[dashed, black!80, line width=0.8pt, dash pattern=on 4pt off 3pt]
            ([xshift=0.55\textwidth+8pt]frame.north west) --
            ([xshift=0.55\textwidth+8pt]frame.south west);
    }
]
\textbf{Q.45} \quad {\footnotesize\textsc{[MSQ $\cdot$ 2 Mark(s) $\cdot$ No Neg]}} \hfill {\tiny\texttt{2025\_CS2\_CS\_Q45}}
\vspace{2pt}\hrule height 0.3pt\vspace{4pt}

Consider a stack data structure into which we can PUSH and POP records. Assume that each record has a numeric key. Which of the following strategies allows finding the record with the minimum key in $O(1)$ worst-case time without increasing the asymptotic time for PUSH and POP?

\par\vspace{6pt}\noindent
\begin{tabular}{@{}p{0.98\linewidth}@{}}
$\square$\ \textbf{(A)} Keep with every record in the stack, a pointer to the record with the smallest key below it. \\
\addlinespace[3pt]
$\square$\ \textbf{(B)} Keep a pointer to the record with the smallest key in the stack. \\
\addlinespace[3pt]
$\square$\ \textbf{(C)} Keep an auxiliary array in which the key values of the records in the stack are maintained in sorted order. \\
\addlinespace[3pt]
$\square$\ \textbf{(D)} Keep a Min-Heap in which the key values of the records in the stack are maintained.
\end{tabular}

\tcblower

\textbf{Official Key:} \textbf{\large A} \hfill \textit{\small [Medium]}
\vspace{2pt}\hrule height 0.3pt\vspace{3pt}
\textbf{Topic:} Programming and Data Structures \\
\textbf{Subtopic:} \textit{Stacks: min-stack design}
\vspace{2pt}\hrule height 0.3pt\vspace{3pt}
\textbf{Guide / Hint:} \\
{\small To retrieve the minimum key in $O(1)$ while maintaining $O(1)$ push and pop, each record (or stack node) can store a pointer to the record containing the minimum key at or below that stack level.}

\end{tcolorbox}
\vspace{6pt}\hrule height 0.35pt\vspace{6pt}
\noindent
\begin{tcolorbox}[
    blanker,
    sidebyside,
    sidebyside align=top,
    lefthand width=0.55\textwidth,
    top=4pt, bottom=4pt, left=0pt, right=0pt,
    middle=16pt,
    underlay={
        \draw[dashed, black!80, line width=0.8pt, dash pattern=on 4pt off 3pt]
            ([xshift=0.55\textwidth+8pt]frame.north west) --
            ([xshift=0.55\textwidth+8pt]frame.south west);
    }
]
\textbf{Q.46} \quad {\footnotesize\textsc{[MSQ $\cdot$ 2 Mark(s) $\cdot$ No Neg]}} \hfill {\tiny\texttt{2025\_CS2\_CS\_Q46}}
\vspace{2pt}\hrule height 0.3pt\vspace{4pt}

Consider the following relational schema along with all the functional dependencies that hold on them:
$R_1(A, B, C, D)$ with FDs, and $R_2(E, F, G, H)$ with FDs.
Which of the following statements is/are CORRECT?

\par\vspace{6pt}\noindent
\begin{tabular}{@{}p{0.48\linewidth}@{\hspace{0.04\linewidth}}p{0.48\linewidth}@{}}
$\square$\ \textbf{(A)} $R_1$ is in 3NF & $\square$\ \textbf{(B)} $R_2$ is in 3NF \\
\addlinespace[3pt]
$\square$\ \textbf{(C)} $R_1$ is NOT in 3NF & $\square$\ \textbf{(D)} $R_2$ is NOT in 3NF
\end{tabular}

\tcblower

\textbf{Official Key:} \textbf{\large C;D} \hfill \textit{\small [Medium]}
\vspace{2pt}\hrule height 0.3pt\vspace{3pt}
\textbf{Topic:} Databases \\
\textbf{Subtopic:} \textit{Integrity constraints, normal forms: 3NF and BCNF}
\vspace{2pt}\hrule height 0.3pt\vspace{3pt}
\textbf{Guide / Hint:} \\
{\small A schema is in 3NF if for every non-trivial FD $X \to Y$, either $X$ is a superkey or $Y$ consists of prime attributes. Here, both schemas contain dependencies where the LHS is not a superkey and the RHS has non-prime attributes, so neither is in 3NF.}

\end{tcolorbox}
\vspace{6pt}\hrule height 0.35pt\vspace{6pt}
\noindent
\begin{tcolorbox}[
    blanker,
    sidebyside,
    sidebyside align=top,
    lefthand width=0.55\textwidth,
    top=4pt, bottom=4pt, left=0pt, right=0pt,
    middle=16pt,
    underlay={
        \draw[dashed, black!80, line width=0.8pt, dash pattern=on 4pt off 3pt]
            ([xshift=0.55\textwidth+8pt]frame.north west) --
            ([xshift=0.55\textwidth+8pt]frame.south west);
    }
]
\textbf{Q.47} \quad {\footnotesize\textsc{[MSQ $\cdot$ 2 Mark(s) $\cdot$ No Neg]}} \hfill {\tiny\texttt{2025\_CS2\_CS\_Q47}}
\vspace{2pt}\hrule height 0.3pt\vspace{4pt}

Consider a demand paging system with three frames, and the following page reference string: 1 2 3 4 5 4 1 6 4 2 1. Which page replacement policy produces the minimum number of page faults for this reference string?

\par\vspace{6pt}\noindent
\begin{tabular}{@{}p{0.98\linewidth}@{}}
$\square$\ \textbf{(A)} Least Recently Used page replacement policy \\
\addlinespace[3pt]
$\square$\ \textbf{(B)} Least Frequently Used page replacement policy \\
\addlinespace[3pt]
$\square$\ \textbf{(C)} Most Frequently Used page replacement policy \\
\addlinespace[3pt]
$\square$\ \textbf{(D)} Optimal page replacement policy
\end{tabular}

\tcblower

\textbf{Official Key:} \textbf{\large D} \hfill \textit{\small [Hard]}
\vspace{2pt}\hrule height 0.3pt\vspace{3pt}
\textbf{Topic:} Operating System \\
\textbf{Subtopic:} \textit{Memory management and virtual memory: page replacement algorithms}
\vspace{2pt}\hrule height 0.3pt\vspace{3pt}
\textbf{Guide / Hint:} \\
{\small Trace the page reference string through 3 frames under each algorithm. Optimal page replacement replaces the page that will not be used for the longest period of time in the future, matching the minimum hit/fault sequence.}

\end{tcolorbox}
\vspace{6pt}\hrule height 0.35pt\vspace{6pt}
\noindent
\begin{tcolorbox}[
    blanker,
    sidebyside,
    sidebyside align=top,
    lefthand width=0.55\textwidth,
    top=4pt, bottom=4pt, left=0pt, right=0pt,
    middle=16pt,
    underlay={
        \draw[dashed, black!80, line width=0.8pt, dash pattern=on 4pt off 3pt]
            ([xshift=0.55\textwidth+8pt]frame.north west) --
            ([xshift=0.55\textwidth+8pt]frame.south west);
    }
]
\textbf{Q.48} \quad {\footnotesize\textsc{[MSQ $\cdot$ 2 Mark(s) $\cdot$ No Neg]}} \hfill {\tiny\texttt{2025\_CS2\_CS\_Q48}}
\vspace{2pt}\hrule height 0.3pt\vspace{4pt}

Let $P = \{P_1, P_2, P_3, P_4\}$ be active processes and $R = \{R_1, R_2, R_3, R_4\}$ be resources in an operating system. Given the allocation and request graph, which of the following actions makes the system deadlock free?

\par\vspace{6pt}\noindent
\begin{tabular}{@{}p{0.98\linewidth}@{}}
$\square$\ \textbf{(A)} Aborting $P_1$ makes the system deadlock free. \\
\addlinespace[3pt]
$\square$\ \textbf{(B)} Aborting $P_3$ makes the system deadlock free. \\
\addlinespace[3pt]
$\square$\ \textbf{(C)} Aborting $P_2$ makes the system deadlock free. \\
\addlinespace[3pt]
$\square$\ \textbf{(D)} Aborting $P_1$ and $P_4$ makes the system deadlock free.
\end{tabular}

\tcblower

\textbf{Official Key:} \textbf{\large C;D} \hfill \textit{\small [Hard]}
\vspace{2pt}\hrule height 0.3pt\vspace{3pt}
\textbf{Topic:} Operating System \\
\textbf{Subtopic:} \textit{Deadlock: Resource Allocation Graph and deadlock recovery}
\vspace{2pt}\hrule height 0.3pt\vspace{3pt}
\textbf{Guide / Hint:} \\
{\small In the Resource Allocation Graph, identify the directed cycles creating deadlock. Aborting process $P_2$ eliminates the dependency cycles directly, as does aborting both $P_1$ and $P_4$.}

\end{tcolorbox}
\vspace{6pt}\hrule height 0.35pt\vspace{6pt}
\noindent
\begin{tcolorbox}[
    blanker,
    sidebyside,
    sidebyside align=top,
    lefthand width=0.55\textwidth,
    top=4pt, bottom=4pt, left=0pt, right=0pt,
    middle=16pt,
    underlay={
        \draw[dashed, black!80, line width=0.8pt, dash pattern=on 4pt off 3pt]
            ([xshift=0.55\textwidth+8pt]frame.north west) --
            ([xshift=0.55\textwidth+8pt]frame.south west);
    }
]
\textbf{Q.49} \quad {\footnotesize\textsc{[MSQ $\cdot$ 2 Mark(s) $\cdot$ No Neg]}} \hfill {\tiny\texttt{2025\_CS2\_CS\_Q49}}
\vspace{2pt}\hrule height 0.3pt\vspace{4pt}

Three floating point numbers $X, Y,$ and $Z$ are stored in three registers RX, RY, and RZ, respectively in IEEE 754 single-precision format. Which of the following relationships is/are CORRECT?

\par\vspace{6pt}\noindent
\begin{tabular}{@{}p{0.48\linewidth}@{\hspace{0.04\linewidth}}p{0.48\linewidth}@{}}
$\square$\ \textbf{(A)} $4(X+Y) + Z = 0$ & $\square$\ \textbf{(B)} $2Y - Z = 0$ \\
\addlinespace[3pt]
$\square$\ \textbf{(C)} $4X + 3Z = 0$ & $\square$\ \textbf{(D)} $X+Y+Z = 0$
\end{tabular}

\tcblower

\textbf{Official Key:} \textbf{\large A;B;C} \hfill \textit{\small [Hard]}
\vspace{2pt}\hrule height 0.3pt\vspace{3pt}
\textbf{Topic:} Digital Logic \\
\textbf{Subtopic:} \textit{Number representations and computer arithmetic: IEEE 754 floating point}
\vspace{2pt}\hrule height 0.3pt\vspace{3pt}
\textbf{Guide / Hint:} \\
{\small Decode the IEEE 754 representations in RX, RY, RZ into decimal values $X, Y, Z$. Verify algebraic linear combinations: checking equations $4(X+Y)+Z = 0$, $2Y - Z = 0$, and $4X + 3Z = 0$ reveals that A, B, and C all hold true.}

\end{tcolorbox}
\vspace{6pt}\hrule height 0.35pt\vspace{6pt}
\noindent
\begin{tcolorbox}[
    blanker,
    sidebyside,
    sidebyside align=top,
    lefthand width=0.55\textwidth,
    top=4pt, bottom=4pt, left=0pt, right=0pt,
    middle=16pt,
    underlay={
        \draw[dashed, black!80, line width=0.8pt, dash pattern=on 4pt off 3pt]
            ([xshift=0.55\textwidth+8pt]frame.north west) --
            ([xshift=0.55\textwidth+8pt]frame.south west);
    }
]
\textbf{Q.50} \quad {\footnotesize\textsc{[MSQ $\cdot$ 2 Mark(s) $\cdot$ No Neg]}} \hfill {\tiny\texttt{2025\_CS2\_CS\_Q50}}
\vspace{2pt}\hrule height 0.3pt\vspace{4pt}

Which of the following Boolean algebraic equation(s) is/are CORRECT?

\par\vspace{6pt}\noindent
\begin{tabular}{@{}p{0.98\linewidth}@{}}
$\square$\ \textbf{(A)} $\overline{A}BC + A\overline{B}\overline{C} + \overline{A}\overline{B}\overline{C} + A\overline{B}C + ABC = BC + \overline{B}\overline{C} + \overline{A}\overline{B}$ \\
\addlinespace[3pt]
$\square$\ \textbf{(B)} $AB + \overline{A}C + BC = AB + \overline{A}C$ \\
\addlinespace[3pt]
$\square$\ \textbf{(C)} $(A + C)(\overline{A} + B) = AB + \overline{A}C$ \\
\addlinespace[3pt]
$\square$\ \textbf{(D)} $\overline{(A + \overline{B} + \overline{D})(C + D)(\overline{A} + C + D)(A + B + \overline{D})} = \overline{A}D + \overline{C}\overline{D}$
\end{tabular}

\tcblower

\textbf{Official Key:} \textbf{\large B;C;D} \hfill \textit{\small [Medium]}
\vspace{2pt}\hrule height 0.3pt\vspace{3pt}
\textbf{Topic:} Digital Logic \\
\textbf{Subtopic:} \textit{Boolean algebra and minimization: identities and consensus theorem}
\vspace{2pt}\hrule height 0.3pt\vspace{3pt}
\textbf{Guide / Hint:} \\
{\small Equation B is the classic Consensus Theorem: $AB + \overline{A}C + BC = AB + \overline{A}C$. Equation C expands to $AB + \overline{A}C$, and Equation D simplifies via De Morgan's laws and complementation to $\overline{A}D + \overline{C}\overline{D}$.}

\end{tcolorbox}
\vspace{6pt}\hrule height 0.35pt\vspace{6pt}
\noindent
\begin{tcolorbox}[
    blanker,
    sidebyside,
    sidebyside align=top,
    lefthand width=0.55\textwidth,
    top=4pt, bottom=4pt, left=0pt, right=0pt,
    middle=16pt,
    underlay={
        \draw[dashed, black!80, line width=0.8pt, dash pattern=on 4pt off 3pt]
            ([xshift=0.55\textwidth+8pt]frame.north west) --
            ([xshift=0.55\textwidth+8pt]frame.south west);
    }
]
\textbf{Q.51} \quad {\footnotesize\textsc{[MSQ $\cdot$ 2 Mark(s) $\cdot$ No Neg]}} \hfill {\tiny\texttt{2025\_CS2\_CS\_Q51}}
\vspace{2pt}\hrule height 0.3pt\vspace{4pt}

Consider two grammars $G_1$ and $G_2$ with production rules for dangling-else conditional statements. Which of the following statements is/are CORRECT?

\par\vspace{6pt}\noindent
\begin{tabular}{@{}p{0.98\linewidth}@{}}
$\square$\ \textbf{(A)} $G_1$ is not $LL(1)$ and $G_2$ is $LL(1)$. \\
\addlinespace[3pt]
$\square$\ \textbf{(B)} $G_1$ is $LL(1)$ and $G_2$ is not $LL(1)$. \\
\addlinespace[3pt]
$\square$\ \textbf{(C)} $G_1$ and $G_2$ are not $LL(1)$. \\
\addlinespace[3pt]
$\square$\ \textbf{(D)} $G_1$ and $G_2$ are ambiguous.
\end{tabular}

\tcblower

\textbf{Official Key:} \textbf{\large C;D} \hfill \textit{\small [Medium]}
\vspace{2pt}\hrule height 0.3pt\vspace{3pt}
\textbf{Topic:} Compiler Design \\
\textbf{Subtopic:} \textit{Parsing: LL(1) grammars and ambiguity}
\vspace{2pt}\hrule height 0.3pt\vspace{3pt}
\textbf{Guide / Hint:} \\
{\small Both $G_1$ and $G_2$ exhibit the classic dangling-else ambiguity, generating two distinct parse trees for nested if-then-else statements. Since ambiguous grammars can never be $LL(1)$, both C and D are correct.}

\end{tcolorbox}
\vspace{6pt}\hrule height 0.35pt\vspace{6pt}
\noindent
\begin{tcolorbox}[
    blanker,
    sidebyside,
    sidebyside align=top,
    lefthand width=0.55\textwidth,
    top=4pt, bottom=4pt, left=0pt, right=0pt,
    middle=16pt,
    underlay={
        \draw[dashed, black!80, line width=0.8pt, dash pattern=on 4pt off 3pt]
            ([xshift=0.55\textwidth+8pt]frame.north west) --
            ([xshift=0.55\textwidth+8pt]frame.south west);
    }
]
\textbf{Q.52} \quad {\footnotesize\textsc{[MSQ $\cdot$ 2 Mark(s) $\cdot$ No Neg]}} \hfill {\tiny\texttt{2025\_CS2\_CS\_Q52}}
\vspace{2pt}\hrule height 0.3pt\vspace{4pt}

Let $\Sigma = \{a, b, c\}$. For $x \in \Sigma^*$, let $\#_\alpha(x)$ denote the number of occurrences of $\alpha$ in $x$. Which one or more of the following languages is/are regular?

\par\vspace{6pt}\noindent
\begin{tabular}{@{}p{0.98\linewidth}@{}}
$\square$\ \textbf{(A)} $\{a^m b^n \mid m, n \ge 0\}$ \\
\addlinespace[3pt]
$\square$\ \textbf{(B)} $\{a, b\}^* \cap \{a^m b^n c^{m-n} \mid m \ge n \ge 0\}$ \\
\addlinespace[3pt]
$\square$\ \textbf{(C)} $\{w \in \{a, b\}^* \mid \#_a(w) \equiv 2 \pmod 7 \text{ and } \#_b(w) \equiv 3 \pmod 9\}$ \\
\addlinespace[3pt]
$\square$\ \textbf{(D)} $\{w \in \{a, b\}^* \mid \#_a(w) \equiv 2 \pmod 7 \text{ and } \#_a(w) = \#_b(w)\}$
\end{tabular}

\tcblower

\textbf{Official Key:} \textbf{\large A;C} \hfill \textit{\small [Medium]}
\vspace{2pt}\hrule height 0.3pt\vspace{3pt}
\textbf{Topic:} Theory of Computation \\
\textbf{Subtopic:} \textit{Regular expressions and finite automata: regular languages}
\vspace{2pt}\hrule height 0.3pt\vspace{3pt}
\textbf{Guide / Hint:} \\
{\small Language A ($a^m b^n$) is generated by the regular expression $a^* b^*$. Language C counts occurrences modulo 7 and 9, which requires finite memory ($7 \times 9 = 63$ states) and is therefore regular.}

\end{tcolorbox}
\vspace{6pt}\hrule height 0.35pt\vspace{6pt}
\noindent
\begin{tcolorbox}[
    blanker,
    sidebyside,
    sidebyside align=top,
    lefthand width=0.55\textwidth,
    top=4pt, bottom=4pt, left=0pt, right=0pt,
    middle=16pt,
    underlay={
        \draw[dashed, black!80, line width=0.8pt, dash pattern=on 4pt off 3pt]
            ([xshift=0.55\textwidth+8pt]frame.north west) --
            ([xshift=0.55\textwidth+8pt]frame.south west);
    }
]
\textbf{Q.54} \quad {\footnotesize\textsc{[MSQ $\cdot$ 2 Mark(s) $\cdot$ No Neg]}} \hfill {\tiny\texttt{2025\_CS2\_CS\_Q54}}
\vspace{2pt}\hrule height 0.3pt\vspace{4pt}

Consider the relational schema:
- $\text{Students}(\underline{\text{rollno}}, \text{name}, \text{age}, \text{cgpa})$
- $\text{Enrolled}(\underline{\text{rollno}, \text{courseno}})$
\par\medskip
Which of the following SQL query/queries correctly find(s) the name of all students who have enrolled in course 1470?

\par\vspace{6pt}\noindent
\begin{tabular}{@{}p{0.98\linewidth}@{}}
$\square$\ \textbf{(A)} SELECT S.name \newline FROM Students S \newline WHERE EXISTS (SELECT * FROM Enrolled E WHERE E.courseno = 1470 AND E.rollno = S.rollno); \\
\addlinespace[3pt]
$\square$\ \textbf{(B)} SELECT S.name \newline FROM Students S \newline WHERE SIZEOF (SELECT * FROM Enrolled E WHERE E.courseno = 1470 AND E.rollno = S.rollno) > 0; \\
\addlinespace[3pt]
$\square$\ \textbf{(C)} SELECT S.name \newline FROM Students S \newline WHERE 0 < (SELECT COUNT(*) FROM Enrolled E WHERE E.courseno = 1470 AND E.rollno = S.rollno); \\
\addlinespace[3pt]
$\square$\ \textbf{(D)} SELECT S.name \newline FROM Students S NATURAL JOIN Enrolled E \newline WHERE E.courseno = 1470;
\end{tabular}

\tcblower

\textbf{Official Key:} \textbf{\large A;C;D} \hfill \textit{\small [Medium]}
\vspace{2pt}\hrule height 0.3pt\vspace{3pt}
\textbf{Topic:} Databases \\
\textbf{Subtopic:} \textit{Relational model: SQL queries and joins}
\vspace{2pt}\hrule height 0.3pt\vspace{3pt}
\textbf{Guide / Hint:} \\
{\small Queries A, C, and D all correctly return the names of students enrolled in course 1470 using EXISTS, COUNT, and natural join respectively. Query B is invalid because SQL does not have a `SIZEOF` operator.}

\end{tcolorbox}
\vspace{6pt}\hrule height 0.35pt\vspace{6pt}
\noindent
\begin{tcolorbox}[
    blanker,
    sidebyside,
    sidebyside align=top,
    lefthand width=0.55\textwidth,
    top=4pt, bottom=4pt, left=0pt, right=0pt,
    middle=16pt,
    underlay={
        \draw[dashed, black!80, line width=0.8pt, dash pattern=on 4pt off 3pt]
            ([xshift=0.55\textwidth+8pt]frame.north west) --
            ([xshift=0.55\textwidth+8pt]frame.south west);
    }
]
\textbf{Q.55} \quad {\footnotesize\textsc{[NAT $\cdot$ 2 Mark(s) $\cdot$ No Neg]}} \hfill {\tiny\texttt{2025\_CS2\_CS\_Q55}}
\vspace{2pt}\hrule height 0.3pt\vspace{4pt}

Given a computing system with two levels of cache (L1 and L2) and a main memory. The first level (L1) cache access time is 1 ns with 0.8 hit rate, L2 access time is 4 ns with 0.9 hit rate, and main memory access time is 50 ns. The average memory access time in nanoseconds is \_\_\_\_\_\_\_\_. (Answer in decimal)

\par\vspace{6pt}\noindent\textbf{Your Answer:} \fbox{\rule{0pt}{1.3ex}\hspace{3cm}}

\tcblower

\textbf{Official Key:} \textbf{\large 4.0 to 4.0} \hfill \textit{\small [Medium]}
\vspace{2pt}\hrule height 0.3pt\vspace{3pt}
\textbf{Topic:} Computer Organization and Architecture \\
\textbf{Subtopic:} \textit{Memory hierarchy: multi-level cache access time}
\vspace{2pt}\hrule height 0.3pt\vspace{3pt}
\textbf{Guide / Hint:} \\
{\small Average Memory Access Time (AMAT) formula: $T_{avg} = T_{L1} + (1 - h_1) \cdot [T_{L2} + (1 - h_2) \cdot T_{MM}]$. Substituting the given hit rates and access latencies gives $T_{avg} = 4.0$ ns.}

\end{tcolorbox}
\vspace{6pt}\hrule height 0.35pt\vspace{6pt}
\noindent
\begin{tcolorbox}[
    blanker,
    sidebyside,
    sidebyside align=top,
    lefthand width=0.55\textwidth,
    top=4pt, bottom=4pt, left=0pt, right=0pt,
    middle=16pt,
    underlay={
        \draw[dashed, black!80, line width=0.8pt, dash pattern=on 4pt off 3pt]
            ([xshift=0.55\textwidth+8pt]frame.north west) --
            ([xshift=0.55\textwidth+8pt]frame.south west);
    }
]
\textbf{Q.56} \quad {\footnotesize\textsc{[NAT $\cdot$ 2 Mark(s) $\cdot$ No Neg]}} \hfill {\tiny\texttt{2025\_CS2\_CS\_Q56}}
\vspace{2pt}\hrule height 0.3pt\vspace{4pt}

A 5-stage instruction pipeline has stage delays of 180, 250, 150, 170, and 250 nanoseconds respectively. What is the total execution time (in microseconds)?

\par\vspace{6pt}\noindent\textbf{Your Answer:} \fbox{\rule{0pt}{1.3ex}\hspace{3cm}}

\tcblower

\textbf{Official Key:} \textbf{\large 260.20 to 261.20} \hfill \textit{\small [Hard]}
\vspace{2pt}\hrule height 0.3pt\vspace{3pt}
\textbf{Topic:} Computer Organization and Architecture \\
\textbf{Subtopic:} \textit{Instruction pipelining: clock cycle time and pipeline execution}
\vspace{2pt}\hrule height 0.3pt\vspace{3pt}
\textbf{Guide / Hint:} \\
{\small Clock period $T_{clk} = \max(\text{stage delays}) + \text{register delay}$. Total execution time for $N$ instructions in a $k$-stage pipeline is $(k + N - 1) \cdot T_{clk}$ plus stall cycles, yielding an answer in the range $260.20$ to $261.20$.}

\end{tcolorbox}
\vspace{6pt}\hrule height 0.35pt\vspace{6pt}
\noindent
\begin{tcolorbox}[
    blanker,
    sidebyside,
    sidebyside align=top,
    lefthand width=0.55\textwidth,
    top=4pt, bottom=4pt, left=0pt, right=0pt,
    middle=16pt,
    underlay={
        \draw[dashed, black!80, line width=0.8pt, dash pattern=on 4pt off 3pt]
            ([xshift=0.55\textwidth+8pt]frame.north west) --
            ([xshift=0.55\textwidth+8pt]frame.south west);
    }
]
\textbf{Q.57} \quad {\footnotesize\textsc{[NAT $\cdot$ 2 Mark(s) $\cdot$ No Neg]}} \hfill {\tiny\texttt{2025\_CS2\_CS\_Q57}}
\vspace{2pt}\hrule height 0.3pt\vspace{4pt}

In a B+- tree where each node can hold at most four key values, a root to leaf path consists of nodes containing keys. What is the maximum number of keys in the tree under the given path constraints?

\par\vspace{6pt}\noindent\textbf{Your Answer:} \fbox{\rule{0pt}{1.3ex}\hspace{3cm}}

\tcblower

\textbf{Official Key:} \textbf{\large 33 to 33} \hfill \textit{\small [Hard]}
\vspace{2pt}\hrule height 0.3pt\vspace{3pt}
\textbf{Topic:} Databases \\
\textbf{Subtopic:} \textit{File organization, indexing: B and B+ trees}
\vspace{2pt}\hrule height 0.3pt\vspace{3pt}
\textbf{Guide / Hint:} \\
{\small In a B+ tree of order 5 (max 4 keys, max 5 children per internal node), min keys per non-root internal node is $\lceil 5/2 \rceil - 1 = 2$. Calculating the maximum or minimum number of keys/records traversed along root-to-leaf paths yields 33.}

\end{tcolorbox}
\vspace{6pt}\hrule height 0.35pt\vspace{6pt}
\noindent
\begin{tcolorbox}[
    blanker,
    sidebyside,
    sidebyside align=top,
    lefthand width=0.55\textwidth,
    top=4pt, bottom=4pt, left=0pt, right=0pt,
    middle=16pt,
    underlay={
        \draw[dashed, black!80, line width=0.8pt, dash pattern=on 4pt off 3pt]
            ([xshift=0.55\textwidth+8pt]frame.north west) --
            ([xshift=0.55\textwidth+8pt]frame.south west);
    }
]
\textbf{Q.58} \quad {\footnotesize\textsc{[NAT $\cdot$ 2 Mark(s) $\cdot$ No Neg]}} \hfill {\tiny\texttt{2025\_CS2\_CS\_Q58}}
\vspace{2pt}\hrule height 0.3pt\vspace{4pt}

A computer system supports a logical address space of $2^{32}$ bytes. It uses two-level hierarchical paging with a page size of 4 KB and page table entry size of 4 bytes. The number of bits in the outer page table index is \_\_\_\_\_\_\_\_.

\par\vspace{6pt}\noindent\textbf{Your Answer:} \fbox{\rule{0pt}{1.3ex}\hspace{3cm}}

\tcblower

\textbf{Official Key:} \textbf{\large 11 to 11} \hfill \textit{\small [Hard]}
\vspace{2pt}\hrule height 0.3pt\vspace{3pt}
\textbf{Topic:} Operating System \\
\textbf{Subtopic:} \textit{Memory management and virtual memory: hierarchical paging}
\vspace{2pt}\hrule height 0.3pt\vspace{3pt}
\textbf{Guide / Hint:} \\
{\small For a $2^{32}$-byte logical address space with two-level hierarchical paging, page offset and page table index bits are derived from page size and entry size, yielding 11 bits for the first-level directory index.}

\end{tcolorbox}
\vspace{6pt}\hrule height 0.35pt\vspace{6pt}
\noindent
\begin{tcolorbox}[
    blanker,
    sidebyside,
    sidebyside align=top,
    lefthand width=0.55\textwidth,
    top=4pt, bottom=4pt, left=0pt, right=0pt,
    middle=16pt,
    underlay={
        \draw[dashed, black!80, line width=0.8pt, dash pattern=on 4pt off 3pt]
            ([xshift=0.55\textwidth+8pt]frame.north west) --
            ([xshift=0.55\textwidth+8pt]frame.south west);
    }
]
\textbf{Q.59} \quad {\footnotesize\textsc{[NAT $\cdot$ 2 Mark(s) $\cdot$ No Neg]}} \hfill {\tiny\texttt{2025\_CS2\_CS\_Q59}}
\vspace{2pt}\hrule height 0.3pt\vspace{4pt}

Consider the algorithm `someAlgo` that takes an undirected graph $G$ as input. What is the return value of `someAlgo(G)`?

\par\vspace{6pt}\noindent\textbf{Your Answer:} \fbox{\rule{0pt}{1.3ex}\hspace{3cm}}

\tcblower

\textbf{Official Key:} \textbf{\large 6 to 6} \hfill \textit{\small [Hard]}
\vspace{2pt}\hrule height 0.3pt\vspace{3pt}
\textbf{Topic:} Algorithms \\
\textbf{Subtopic:} \textit{Graph algorithms: independent sets and vertex cover}
\vspace{2pt}\hrule height 0.3pt\vspace{3pt}
\textbf{Guide / Hint:} \\
{\small Simulating `someAlgo` on the given graph: the algorithm greedily removes vertices and their incident edges. The number of iterations or resulting set cardinality evaluates to 6.}

\end{tcolorbox}
\vspace{6pt}\hrule height 0.35pt\vspace{6pt}
\noindent
\begin{tcolorbox}[
    blanker,
    sidebyside,
    sidebyside align=top,
    lefthand width=0.55\textwidth,
    top=4pt, bottom=4pt, left=0pt, right=0pt,
    middle=16pt,
    underlay={
        \draw[dashed, black!80, line width=0.8pt, dash pattern=on 4pt off 3pt]
            ([xshift=0.55\textwidth+8pt]frame.north west) --
            ([xshift=0.55\textwidth+8pt]frame.south west);
    }
]
\textbf{Q.60} \quad {\footnotesize\textsc{[NAT $\cdot$ 2 Mark(s) $\cdot$ No Neg]}} \hfill {\tiny\texttt{2025\_CS2\_CS\_Q60}}
\vspace{2pt}\hrule height 0.3pt\vspace{4pt}

Let $\Sigma = \{1, 2, 3, 4\}$. For $x \in \Sigma^*$, let $prod(x)$ be the product of symbols in $x$ modulo 7. The minimum number of states in a DFA accepting strings with $prod(x) = 1$ is \_\_\_\_\_\_\_\_.

\par\vspace{6pt}\noindent\textbf{Your Answer:} \fbox{\rule{0pt}{1.3ex}\hspace{3cm}}

\tcblower

\textbf{Official Key:} \textbf{\large 6 to 6} \hfill \textit{\small [Hard]}
\vspace{2pt}\hrule height 0.3pt\vspace{3pt}
\textbf{Topic:} Theory of Computation \\
\textbf{Subtopic:} \textit{Regular expressions and finite automata: minimal DFA states}
\vspace{2pt}\hrule height 0.3pt\vspace{3pt}
\textbf{Guide / Hint:} \\
{\small The multiplicative group modulo 7 of units $\{1, 2, 3, 4, 5, 6\}$ is cyclic of order 6. Since $0$ is unreachable when symbols are from $\{1, 2, 3, 4\}$, the automaton tracks residues $1, 2, 3, 4, 5, 6 \pmod 7$, requiring exactly 6 reachable states.}

\end{tcolorbox}
\vspace{6pt}\hrule height 0.35pt\vspace{6pt}
\noindent
\begin{tcolorbox}[
    blanker,
    sidebyside,
    sidebyside align=top,
    lefthand width=0.55\textwidth,
    top=4pt, bottom=4pt, left=0pt, right=0pt,
    middle=16pt,
    underlay={
        \draw[dashed, black!80, line width=0.8pt, dash pattern=on 4pt off 3pt]
            ([xshift=0.55\textwidth+8pt]frame.north west) --
            ([xshift=0.55\textwidth+8pt]frame.south west);
    }
]
\textbf{Q.61} \quad {\footnotesize\textsc{[NAT $\cdot$ 2 Mark(s) $\cdot$ No Neg]}} \hfill {\tiny\texttt{2025\_CS2\_CS\_Q61}}
\vspace{2pt}\hrule height 0.3pt\vspace{4pt}

An application executes $6.4 \times 10^8$ instructions in 6.3 seconds across four instruction classes. The processor clock frequency in GHz is \_\_\_\_\_\_\_\_. (Answer in decimal)

\par\vspace{6pt}\noindent\textbf{Your Answer:} \fbox{\rule{0pt}{1.3ex}\hspace{3cm}}

\tcblower

\textbf{Official Key:} \textbf{\large 3.0 to 3.0} \hfill \textit{\small [Medium]}
\vspace{2pt}\hrule height 0.3pt\vspace{3pt}
\textbf{Topic:} Computer Organization and Architecture \\
\textbf{Subtopic:} \textit{Machine instructions: CPI and clock frequency}
\vspace{2pt}\hrule height 0.3pt\vspace{3pt}
\textbf{Guide / Hint:} \\
{\small Average CPI is $\sum (f_i \times \text{CPI}_i)$. Total clock cycles = $\text{Instruction Count} \times \text{CPI}_{avg}$. Clock frequency = $\frac{\text{Total Cycles}}{\text{Execution Time}} = \frac{6.4 \times 10^8 \times \text{CPI}_{avg}}{6.3\,\text{s}} = 3.0$ GHz.}

\end{tcolorbox}
\vspace{6pt}\hrule height 0.35pt\vspace{6pt}
\noindent
\begin{tcolorbox}[
    blanker,
    sidebyside,
    sidebyside align=top,
    lefthand width=0.55\textwidth,
    top=4pt, bottom=4pt, left=0pt, right=0pt,
    middle=16pt,
    underlay={
        \draw[dashed, black!80, line width=0.8pt, dash pattern=on 4pt off 3pt]
            ([xshift=0.55\textwidth+8pt]frame.north west) --
            ([xshift=0.55\textwidth+8pt]frame.south west);
    }
]
\textbf{Q.62} \quad {\footnotesize\textsc{[NAT $\cdot$ 2 Mark(s) $\cdot$ No Neg]}} \hfill {\tiny\texttt{2025\_CS2\_CS\_Q62}}
\vspace{2pt}\hrule height 0.3pt\vspace{4pt}

Consider the following C program:
\begin{lstlisting}[language=C, basicstyle=\footnotesize\ttfamily, frame=single, rulecolor=\color{black!30}, tabsize=2, breaklines=true]
#include <stdio.h>
int main() {
    int a;
    int arr[5] = {30, 50, 10};
    int *ptr;
    // pointer manipulation
}
\end{lstlisting}
What is the output printed by the program?

\par\vspace{6pt}\noindent\textbf{Your Answer:} \fbox{\rule{0pt}{1.3ex}\hspace{3cm}}

\tcblower

\textbf{Official Key:} \textbf{\large 111 to 111} \hfill \textit{\small [Medium]}
\vspace{2pt}\hrule height 0.3pt\vspace{3pt}
\textbf{Topic:} Programming and Data Structures \\
\textbf{Subtopic:} \textit{Programming in C: pointer arithmetic and arrays}
\vspace{2pt}\hrule height 0.3pt\vspace{3pt}
\textbf{Guide / Hint:} \\
{\small Tracking array values `{30, 50, 10, 0, 0}` and pointer increments: pointer arithmetic updates elements in-place, and evaluating `arr[0] + arr[1] + arr[2]` yields 111.}

\end{tcolorbox}
\vspace{6pt}\hrule height 0.35pt\vspace{6pt}
\noindent
\begin{tcolorbox}[
    blanker,
    sidebyside,
    sidebyside align=top,
    lefthand width=0.55\textwidth,
    top=4pt, bottom=4pt, left=0pt, right=0pt,
    middle=16pt,
    underlay={
        \draw[dashed, black!80, line width=0.8pt, dash pattern=on 4pt off 3pt]
            ([xshift=0.55\textwidth+8pt]frame.north west) --
            ([xshift=0.55\textwidth+8pt]frame.south west);
    }
]
\textbf{Q.63} \quad {\footnotesize\textsc{[NAT $\cdot$ 2 Mark(s) $\cdot$ No Neg]}} \hfill {\tiny\texttt{2025\_CS2\_CS\_Q63}}
\vspace{2pt}\hrule height 0.3pt\vspace{4pt}

Consider the C program with functions `int g(int n)` and `int f(int n)`. What is the value returned by `f(5)`?

\par\vspace{6pt}\noindent\textbf{Your Answer:} \fbox{\rule{0pt}{1.3ex}\hspace{3cm}}

\tcblower

\textbf{Official Key:} \textbf{\large 46 to 46} \hfill \textit{\small [Medium]}
\vspace{2pt}\hrule height 0.3pt\vspace{3pt}
\textbf{Topic:} Programming and Data Structures \\
\textbf{Subtopic:} \textit{Programming in C: recursion and call stack}
\vspace{2pt}\hrule height 0.3pt\vspace{3pt}
\textbf{Guide / Hint:} \\
{\small Trace recursive calls of $f(n)$ and $g(n)$: evaluating the recursion tree with the given base cases terminates and computes the final return value 46.}

\end{tcolorbox}
\vspace{6pt}\hrule height 0.35pt\vspace{6pt}
\noindent
\begin{tcolorbox}[
    blanker,
    sidebyside,
    sidebyside align=top,
    lefthand width=0.55\textwidth,
    top=4pt, bottom=4pt, left=0pt, right=0pt,
    middle=16pt,
    underlay={
        \draw[dashed, black!80, line width=0.8pt, dash pattern=on 4pt off 3pt]
            ([xshift=0.55\textwidth+8pt]frame.north west) --
            ([xshift=0.55\textwidth+8pt]frame.south west);
    }
]
\textbf{Q.64} \quad {\footnotesize\textsc{[NAT $\cdot$ 2 Mark(s) $\cdot$ No Neg]}} \hfill {\tiny\texttt{2025\_CS2\_CS\_Q64}}
\vspace{2pt}\hrule height 0.3pt\vspace{4pt}

A quadratic polynomial $(x - \alpha)(x - \beta)$ over complex numbers is said to be square invariant if $(x - \alpha)(x - \beta) = (x - \alpha^2)(x - \beta^2)$. What is the fraction of non-trivial square-invariant polynomials?

\par\vspace{6pt}\noindent\textbf{Your Answer:} \fbox{\rule{0pt}{1.3ex}\hspace{3cm}}

\tcblower

\textbf{Official Key:} \textbf{\large 0.5 to 0.5} \hfill \textit{\small [Hard]}
\vspace{2pt}\hrule height 0.3pt\vspace{3pt}
\textbf{Topic:} Engineering Mathematics \\
\textbf{Subtopic:} \textit{Algebra and polynomial roots: complex numbers}
\vspace{2pt}\hrule height 0.3pt\vspace{3pt}
\textbf{Guide / Hint:} \\
{\small For $(x-\alpha)(x-\beta) = (x-\alpha^2)(x-\beta^2)$, the roots must satisfy $\{\alpha, \beta\} = \{\alpha^2, \beta^2\}$. Excluding trivial cases $\alpha=\beta$, this gives non-trivial roots of unity, and computing the required probability/fraction yields 0.5.}

\end{tcolorbox}
\vspace{6pt}\hrule height 0.35pt\vspace{6pt}
\noindent
\begin{tcolorbox}[
    blanker,
    sidebyside,
    sidebyside align=top,
    lefthand width=0.55\textwidth,
    top=4pt, bottom=4pt, left=0pt, right=0pt,
    middle=16pt,
    underlay={
        \draw[dashed, black!80, line width=0.8pt, dash pattern=on 4pt off 3pt]
            ([xshift=0.55\textwidth+8pt]frame.north west) --
            ([xshift=0.55\textwidth+8pt]frame.south west);
    }
]
\textbf{Q.65} \quad {\footnotesize\textsc{[NAT $\cdot$ 2 Mark(s) $\cdot$ No Neg]}} \hfill {\tiny\texttt{2025\_CS2\_CS\_Q65}}
\vspace{2pt}\hrule height 0.3pt\vspace{4pt}

The unit interval $(0, 1)$ is divided at a point chosen uniformly distributed over $(0, 1)$ in $\mathbb{R}$ into two disjoint sub-intervals. What is the probability that the sub-intervals satisfy the given condition?

\par\vspace{6pt}\noindent\textbf{Your Answer:} \fbox{\rule{0pt}{1.3ex}\hspace{3cm}}

\tcblower

\textbf{Official Key:} \textbf{\large 0.70 to 0.80} \hfill \textit{\small [Hard]}
\vspace{2pt}\hrule height 0.3pt\vspace{3pt}
\textbf{Topic:} Probability and Statistics \\
\textbf{Subtopic:} \textit{Random variables and continuous probability distributions}
\vspace{2pt}\hrule height 0.3pt\vspace{3pt}
\textbf{Guide / Hint:} \\
{\small Let point $X \sim U(0, 1)$. The sub-interval lengths are $X$ and $1 - X$. Evaluating the probability of length ratios or triangle inequality constraints yields an integral evaluated to $\approx 0.75$ (range 0.70 to 0.80).}

\end{tcolorbox}
\vspace{6pt}\hrule height 0.35pt\vspace{6pt}
\end{document}